\chapter{The abelian logarithm}
This chapter constructs the analytic core of the method: the auxiliary class $\tau$, the
translated formal group law $F_\tau$ on the chart at $\tau$, the logarithm $\Lam$ on the
kernel of reduction, the differentials $\omega'_i$ with their regularity, and the pairing
$\lambda$ with its integral representation on residue disks. Fix $q = p^N$ as produced by
Lemma~\ref{lem:tau-exists} and put $\Oq = W(\Fq)$.

\begin{definition}[Regular differentials]
\label{def:regdiff}
\lean{ChabautyColeman.RegularDiff}
\leanok
$\RegDiff := \bigoplus_{i=0}^{g-1} \Qp\cdot x^i\,dx/y$, a $g$-dimensional $\Qp$-vector
space with its $\Zp$-lattice $\RegDiff_{\ZZ} := \bigoplus \Zp\, x^i\,dx/y$. (In the
formalization, differentials are represented by their numerators $h$ with
$\omega = h\,dx/y$, $\deg h < g$; the packaged submodule is formalized in
\texttt{HyperellipticCurves.Archive.RegularDiff}, and the main theorems do not depend on it.)
\end{definition}

\begin{definition}[Good position]
\label{def:goodpos}
\lean{ChabautyColeman.GoodPosition, ChabautyColeman.GoodPosition.isCoprime_right,
  ChabautyColeman.GoodPosition.isCoprime_two_mul_right}
\leanok
\uses{lem:mumford-exists}
A reduced Mumford datum $(u, v)$ over $\Oq$ is in \emph{good position} if $\deg u = g$,
$\bar u$ is squarefree, and $\gcd(\bar u, \bar f) = 1$; the auxiliary data derived from
$\tau$ below are required to satisfy the analogous conditions where they are used. Each
condition is the non-vanishing modulo $p$ of an explicit resultant or discriminant of the
data. The coprimality conditions involving $\bar v$ --- in particular
$\gcd(\bar u, 2\bar v) = 1$, the invertibility used by the chart --- are consequences.
\end{definition}

\begin{lemma}[Abundance of points]
\label{lem:many-points}
\lean{ChabautyColeman.residue_curve_points_infinite,
  ChabautyColeman.infinite_setOf_eval_ne_zero}
\leanok
$\bar C(\Fpbar)$ is infinite; more precisely, the set of non-Weierstrass affine points is
infinite. Hence for any finite set of conditions as in Definition~\ref{def:goodpos} there
are $g$ distinct non-Weierstrass points avoiding them. (Rationality of the chosen points
over a finite field, needed for the counting argument, is handled at the arithmetic
instantiation in Chapter~5.)
\end{lemma}
\begin{proof}
\leanok
For all but finitely many $\bar a \in \Fpbar$, $\bar f(\bar a)$ is a nonzero square in
$\Fpbar$, and each such $\bar a$ lifts to a non-Weierstrass point
$(\bar a, \sqrt{\bar f(\bar a)})$; the conditions exclude only finitely many points.
\end{proof}

\begin{lemma}[Genericity of the centre]
\label{lem:tau-generic}
\lean{ChabautyColeman.HasGoodCenter, ChabautyColeman.hasGoodCenter_of_pos_genus,
  ChabautyColeman.hasGoodCenter_genus_two,
  ChabautyColeman.IsFoldCerts, ChabautyColeman.IsFoldAvoid,
  ChabautyColeman.exists_centerFoldWitness, ChabautyColeman.exists_centerFoldPack,
  ChabautyColeman.exists_isFoldCerts_isFoldAvoid,
  ChabautyColeman.cantorCompose_point_eq_disjoint,
  ChabautyColeman.cantorCompose_point_eq_double,
  ChabautyColeman.isCoprime_foldChain_of_resultant_ne_zero,
  ChabautyColeman.curveAvoidNum, ChabautyColeman.eval_curveAvoidNum,
  ChabautyColeman.not_curveVanishing_curveAvoidNum,
  ChabautyColeman.not_curveVanishing_curveCondNum_of_squarefree,
  ChabautyColeman.not_curveVanishing_curveAvoidNum_of_squarefree,
  ChabautyColeman.not_curveVanishing_X_snd, ChabautyColeman.not_curveVanishing_X_sub_C,
  ChabautyColeman.exists_curveTuple_ne_zero, ChabautyColeman.curveVanishing_mul,
  ChabautyColeman.not_curveVanishing_of_aeval_ne_zero}
\leanok
\uses{def:goodpos, lem:many-points, def:cantor-map, def:centre-fold, thm:fold-witness,
  lem:fold-witness-genus-one}
There exists a tuple of $g$ affine points of the reduced curve over $\Fpbar$ ---
pairwise distinct $x$-coordinates avoiding a prescribed finite set $T$ of residue
$x$-values, nonzero $y$-coordinates --- whose interpolated Mumford datum
$(\bar u, \bar v)$, with $\bar u = \prod_i (X - \bar x_i)$ and $\bar v$ the interpolant
through the points, is such that the point-by-point fold of
Definition~\ref{def:centre-fold} computing $\tau + \tau - \tau$ is defined throughout:
every step condition of the $2g$-step fold is nonzero, and every intermediate state
avoids the Weierstrass points, $\gcd(U_j, 2V_j) = 1$, the starting datum included. Good
position of $(\bar u, \bar v)$ is automatic: $\bar u$ is squarefree because the
$\bar x_i$ are distinct, and coprime to $\bar f$ because the $\bar y_i$ are nonzero. The
statement holds for every genus $g \ge 1$. For $g \ge 2$ no squarefreeness of $\bar f$
is needed; for $g = 1$ it is needed: for $\bar f = X^3$ in characteristic $3$ the fold
conditions vanish identically on the curve
(cf.\ Lemma~\ref{lem:fold-witness-genus-one}). The avoidance of the prescribed set $T$
is what the per-disk configuration of Lemma~\ref{lem:global-generic} requires: $\tau$
must miss the centres of the residue disks being counted.
\end{lemma}
\begin{proof}
\leanok
\uses{thm:fold-witness, lem:fold-witness-genus-one}
The fold witness --- Theorem~\ref{thm:fold-witness} for $g \ge 2$ and
Lemma~\ref{lem:fold-witness-genus-one} for $g = 1$ --- provides a tuple at which the
$2g$ step conditions, the nonvanishing of the $y$-coordinates, the $T$-avoidance and the
per-stage Weierstrass avoidance all hold simultaneously: each condition is a nonzero
polynomial condition on tuples of curve points, and finitely many such conditions have a
common witness (Lemma~\ref{lem:many-points} and the avoidance principle below; the
stage coprimalities enter as resultants of the isobaric numerator states of
Lemma~\ref{lem:fold-universal}, refuted on the cuspidal degeneration). The B\'ezout
certificates required by the certificate-passing formulation of Cantor's algorithm are
then constructed over the field: at a disjoint step, coprimality of $U_j$ and
$X - \bar x_j$ is the step condition $U_j(\bar x_j) \neq 0$; at the doubling step,
coprimality of $X - \bar x_0$ and $V_0 + \bar y_0$ follows from $2\bar y_0 \neq 0$. With
these certificates the certified computation coincides with the fold step by step: at a
disjoint step the composed square root is determined by Chinese-remainder uniqueness,
and at the doubling step by uniqueness of the square root modulo the composed support.
\end{proof}

\bigskip
\noindent\textbf{The witness for all genera: the point-by-point fold and its moment
certificates.} The following sequence of statements proves the existence of the witness
tuple for every genus; the entire sequence is formalized, and its main steps were
additionally verified symbolically at small genus by the scripts in \texttt{prototype/}: the
fold on the cuspidal curve (\texttt{symfold.m}), the moment characterization of its states
and the lemmas behind it (\texttt{apic.m}, \texttt{p1lemmas.m}), and the mirror
(\texttt{revfold.m}); every condition of every genus $\le 6$ has been certified numerically
from the moment determinants alone (\texttt{detcond.m}).

\begin{definition}[The point-by-point fold]
\label{def:centre-fold}
\uses{def:cantor-map, def:goodpos}
\lean{ChabautyColeman.foldChain, ChabautyColeman.IsFoldState,
  ChabautyColeman.IsFoldState.step, ChabautyColeman.cuspChain_eq_foldChain,
  ChabautyColeman.foldGenH, ChabautyColeman.foldH_eq_of_eval_eq,
  ChabautyColeman.foldH_eq_of_dvd_double, ChabautyColeman.foldXs,
  ChabautyColeman.foldYs}
\leanok
Let $k$ be a field of characteristic $\neq 2$, let $f \in k[X]$ be monic squarefree of
degree $N = 2g+1$, and let $P_1, \dots, P_g$ be affine non-Weierstrass points of
$y^2 = f(x)$ with pairwise distinct $x$-coordinates; let $\tau = (u, v)$ be the reduced
Mumford datum interpolating them. The \emph{fold} computes $\tau + \tau - \tau$ by $2g$
single-point additions: $D_0 := \tau$ and $D_j := \operatorname{red}(D_{j-1} + P_j)$ for
$1 \le j \le g$ (\emph{phase 1}), then
$D_{g+m} := \operatorname{red}(D_{g+m-1} + \iota P_{g+1-m})$ for $1 \le m \le g$
(\emph{phase 2}: the involuted points are added in \emph{reverse} order, so that
$[D_{g+m}] = [2\tau] - \sum_{i > g-m}[P_i - \infty] = [D_{g-m}]$ and the fold retraces
its own phase-1 states backwards; see Lemma~\ref{lem:fold-mirror}). Each addition
composes to a divisor of degree $g+1$ whose $b$-polynomial has degree $\le g$, and a
single Cantor reduction returns to degree $g$; since $2\bigl((g+1)-1\bigr) < \deg f$,
every reduction lies in the unconditional degree regime, so no degree-genericity
condition arises at any step. The \emph{conditions} of the fold are: at each step, the
disjointness value $U(x_P)$ of the current $u$-polynomial (at the single doubling step
$j = 1$ instead $U_0'(x_{P_1})$ and $V_0(x_{P_1}) + y_{P_1}$). Further avoidance
quantities --- e.g.\ the Weierstrass avoidance $\operatorname{res}(U_j, f) \neq 0$ of an
intermediate state --- are polynomial in the state coefficients and are treated by the
same universal method.
\end{definition}

\begin{lemma}[Mirror symmetry]
\label{lem:fold-mirror}
\uses{def:centre-fold, def:cusp-fold, thm:cusp-moment}
\lean{ChabautyColeman.congruence_of_even_mul, ChabautyColeman.foldNodes,
  ChabautyColeman.prod_foldNodes_pairing,
  ChabautyColeman.genNode_zero_atom_ne_zero,
  ChabautyColeman.genNode_trio_atom_ne_zero,
  ChabautyColeman.genNode_tripled_atom_ne_zero,
  ChabautyColeman.genCuspFold, ChabautyColeman.genFoldH,
  ChabautyColeman.genCuspFold_eval_sq_ne_zero,
  ChabautyColeman.exists_genCuspFold_witness}
\leanok
On the cuspidal model the hyperelliptic involution is $t \mapsto -t$, so phase 2 of the
fold is the phase-1 recursion continued along the extended node sequence
$(t_1, \dots, t_g, -t_g, \dots, -t_1)$. Assume the fold conditions hold up to phase-2
step $m \le g$. Then the state $A_{g+m}$ of the extended chain equals the phase-1 state
$A_{g-m}$ of the \emph{same} chain: the fold retraces its own states backwards.
Consequently every phase-2 condition value is the evaluation of a phase-1 state at a
single node --- the disjointness value of phase-2 step $m+1$ is $U_{g-m}(x_{g-m})$, the
stage-$(g-m)$ support evaluated at its own last-added node --- and phase 2 needs no
separate analysis: on the cusp, the full $2g$-step fold is controlled by the phase-1
theory alone.
\end{lemma}
\begin{proof}
\leanok
The stage-$(g+m)$ alphabet of the extended chain factors as
$\prod_{i > g-m}(z^2 - t_i^2)$ times the stage-$(g-m)$ alphabet of the same chain ---
an \emph{even} polynomial factor with nonzero value at $0$ at generic nodes. The master
congruence of Theorem~\ref{thm:cusp-moment} descends along such even factors, so
$A_{g+m}$ satisfies the stage-$(g-m)$ moment congruence, and the value-extraction clause
of Theorem~\ref{thm:cusp-moment} reads off every phase-2 condition value against the
phase-1 determinants at the generic nodes (the avoidance value and the two mirror
values, extending the stage-$(g-m)$ alphabet at $\mp t_{g-m}$); no second chain and no
state identification is needed. Only at the final stage $2g$ is $A_{2g} = A_0$ read off
by the uniqueness clause, to get the last avoidance value from $A_0(0) \neq 0$. The
formal argument is an interleaved induction: conditions up to step $m$ give the
descended congruence at $m$, which gives the conditions at $m+1$. At the level of the
curve the identity $D_{g+m} = D_{g-m}$ also follows from the class computation
$[D_{g+m}] = [2\tau] - \sum_{i > g-m}[P_i - \infty] = [\tau] + \sum_{i \le g-m}
[P_i - \infty] = [D_{g-m}]$ together with uniqueness of reduced representatives, but the
formalized route needs only the cuspidal case, where the congruence argument replaces
the Jacobian.
\end{proof}

\begin{lemma}[Universal isobaric numerators]
\label{lem:fold-universal}
\uses{def:centre-fold}
\lean{ChabautyColeman.MumfordNum, ChabautyColeman.MumfordNum.step,
  ChabautyColeman.foldNum, ChabautyColeman.foldHNum, ChabautyColeman.foldCondNum,
  ChabautyColeman.foldBase, ChabautyColeman.foldNum_succ, ChabautyColeman.den_foldNum,
  ChabautyColeman.map_foldNum, ChabautyColeman.map_foldHNum,
  ChabautyColeman.map_foldCondNum, ChabautyColeman.map_foldBase,
  ChabautyColeman.IsScaledState, ChabautyColeman.IsScaledState.step,
  ChabautyColeman.isScaledState_foldNum, ChabautyColeman.foldCond,
  ChabautyColeman.foldCondNum_eq_den_mul, ChabautyColeman.foldStep_dvd_disjoint,
  ChabautyColeman.foldStep_dvd_double, ChabautyColeman.IsIsobaric,
  ChabautyColeman.IsIsobaric.divByMonic, ChabautyColeman.IsIsobaricState,
  ChabautyColeman.IsIsobaricState.step, ChabautyColeman.foldWt,
  ChabautyColeman.foldCondWt, ChabautyColeman.isIsobaricState_foldNum,
  ChabautyColeman.isWeightedHomogeneous_foldCondNum,
  ChabautyColeman.isWeightedHomogeneous_foldHNum,
  ChabautyColeman.isIsobaricState_foldBase, ChabautyColeman.univWt,
  ChabautyColeman.univF, ChabautyColeman.univXs, ChabautyColeman.univYs,
  ChabautyColeman.univBase, ChabautyColeman.univCondNum,
  ChabautyColeman.isWeightedHomogeneous_univCondNum, ChabautyColeman.cuspSpec,
  ChabautyColeman.isWeightedHomogeneous_aeval_cuspSpec,
  ChabautyColeman.cuspNum, ChabautyColeman.isWeightedHomogeneous_cuspNum,
  ChabautyColeman.cuspNum_eq_foldCondNum}
\leanok
Give $X$ and $x_i$ weight $2$, $y_i$ weight $N$, and the coefficient $a_r$ of $f$ weight
$2(N-r)$. There are \emph{universal} polynomials
$\mathcal{N}_q \in \mathbb{Z}[x_1, \dots, x_g,\, y_1, \dots, y_g,\, a_0, \dots, a_{N-1}]$,
isobaric for these weights and indexed by the quantities $q$ of the fold (the
coefficients of every intermediate state and every condition value), such that over any
field of characteristic $\neq 2$ in which the fold is defined,
$q = \mathcal{N}_q(\text{inputs})\big/\bigl(2^{b_q} \prod \mathcal{N}_{q'}^{\,k_{q'}}\bigr)$
with the product over \emph{earlier} condition quantities $q'$. Under the cuspidal
specialization $a = 0$, $x_i = t_i^2$, $y_i = t_i^N$, each $\mathcal{N}_q$ becomes a
homogeneous polynomial $P_q \in \mathbb{Z}[t_1, \dots, t_g]$ of degree
$\operatorname{wt}(\mathcal{N}_q)$.
\end{lemma}
\begin{proof}
\leanok
Induction along the fold in cofactor form. Composition divides only by the current
disjointness value (at the doubling step, by $2 y_{P} U'(x_P)$); the reduction step
divides by nothing: the cofactor division is by a \emph{monic} polynomial, and the
normalization is by the leading coefficient $-1$, as $f$ dominates. Clearing the
denominators inductively produces the $\mathcal{N}_q$, with denominators as stated.
Every operation is equivariant under the scaling
$(x, y, a_r) \mapsto (\lambda^2 x, \lambda^N y, \lambda^{2(N-r)} a_r)$ --- the equation
$y^2 - f(x)$ is isobaric of weight $2N$ --- so the numerators may be chosen isobaric.
The cuspidal specialization sends an isobaric monomial with exponents $\beta_i$ (on
$x_i$) and $\gamma_i$ (on $y_i$) to $\pm t^{c}$ with
$\sum_i c_i = \sum_i (2\beta_i + N\gamma_i)$, which equals the weight when no $a_r$ is
present; monomials containing some $a_r$ vanish.
\end{proof}

\begin{definition}[The cuspidal model and the $A$-recursion]
\label{def:cusp-fold}
\uses{def:centre-fold}
\lean{ChabautyColeman.cuspChain, ChabautyColeman.IsCuspState}
\leanok
The \emph{cuspidal model} is $y^2 = X^N$ with rational parametrization $x = t^2$,
$y = t^N$; Cantor's algorithm operates on Mumford pairs verbatim. Let $t_1, \dots, t_g$
be indeterminates over the prime field ($\operatorname{char} \neq 2$),
$P_i := (t_i^2,\, t_i^N)$, and run the fold over $\mathbb{Q}(t_1, \dots, t_g)$ resp.\
$\mathbb{F}_p(t_1, \dots, t_g)$. Along phase 1 define $A_j \in \mathbb{Q}(t)[z]$ by
$A_0 := \prod_i (z - t_i)$ and, at the step adding $P_{j+1}$ with composed pair
$(U_c, V_c)$: $A_c := A_j \cdot (z - t_{j+1})$ and $A_{j+1} := (-1)^g\, C_c(-z)$, where
$C_c := \bigl(z^N - V_c(z^2)\bigr)/A_c(z)$, a polynomial by
Lemma~\ref{lem:cusp-divisibility}. On the cusp the avoidance condition is
$U_j(0) \neq 0$, since $\operatorname{res}(U, X^N) = \pm U(0)^N$.
\end{definition}

\begin{lemma}[$z$-lift of the Mumford condition]
\label{lem:cusp-divisibility}
\uses{def:cusp-fold}
\lean{ChabautyColeman.cuspCompose_dvd, ChabautyColeman.cuspNext_dvd,
  ChabautyColeman.cuspStep_mumford, ChabautyColeman.IsCuspState.step,
  ChabautyColeman.isCuspState_cuspChain}
\leanok
At every stage of phase 1, $A \mid z^N - V(z^2)$ --- for the reduced states
$(A_j, V_j)$ and for the composed pairs $(A_c, V_c)$ alike.
\end{lemma}
\begin{proof}
\leanok
Induction along the fold. Base: $v$ interpolates, so $v(t_i^2) = t_i^N$ and each
$z - t_i$ divides $z^N - v(z^2)$; the factors are coprime (distinct $t_i$). Composition:
$V_c = V_j + U_j H$ with $H$ the composition scalar, so
$z^N - V_c(z^2) = \bigl(z^N - V_j(z^2)\bigr) - U_j(z^2)H$ is divisible by $A_j$
(inductively, and $A_j(z) \mid U_j(z^2)$ by Lemma~\ref{lem:cusp-reduction}); the new
factor divides because $V_c(t_{j+1}^2) = t_{j+1}^N$ by the interpolation property of the
composition, and it is coprime to $A_j$ by the disjointness condition. Reduction:
$V_{j+1} \equiv -V_c \bmod U_{j+1}$ and $A_{j+1}(z) \mid U_{j+1}(z^2)$, so it suffices
that $A_{j+1} \mid z^N + V_c(z^2) = -\bigl((-z)^N - V_c((-z)^2)\bigr) =
-A_c(-z)C_c(-z)$, which holds since $A_{j+1} = \pm C_c(-z)$.
\end{proof}

\begin{lemma}[Reduction in the $A$-picture]
\label{lem:cusp-reduction}
\uses{def:cusp-fold, lem:cusp-divisibility}
\lean{ChabautyColeman.cuspCofactor_monic, ChabautyColeman.cuspCofactor_natDegree,
  ChabautyColeman.cuspNext_monic, ChabautyColeman.cuspNext_natDegree,
  ChabautyColeman.cuspNext_U, ChabautyColeman.cuspCompose_U_comp}
\leanok
Every $A_j$ is monic of degree $g$, and $U_j(z^2) = (-1)^g A_j(z) A_j(-z)$.
\end{lemma}
\begin{proof}
\leanok
Base: $u(z^2) = \prod(z^2 - t_i^2) = A_0(z) \cdot \prod(z + t_i) =
(-1)^g A_0(z) A_0(-z)$. Step: from $V_c(z^2)^2 - z^{2N} = U_c(z^2)\, w_c(z^2)$ and the
factorizations $V_c(z^2) - z^N = -A_c(z)C_c(z)$ and $V_c(z^2) + z^N = -A_c(-z)C_c(-z)$,
together with $U_c(z^2) = (-1)^{g+1} A_c(z) A_c(-z)$ (composition multiplies $U_j$ by
$X - t_{j+1}^2$), one gets $w_c(z^2) = (-1)^{g+1} C_c(z) C_c(-z)$. The Cantor reduction
returns $U_{j+1} = w_c/\operatorname{lc}(w_c) = -w_c$ (here $\deg w_c = g$ and
$\operatorname{lc}(w_c) = -1$ because $f = X^N$ dominates $V_c^2$), so
$U_{j+1}(z^2) = (-1)^g C_c(z)C_c(-z) = (-1)^g A_{j+1}(z) A_{j+1}(-z)$; and $C_c$ is
monic of degree $g$ ($z^N - V_c(z^2)$ is monic of degree $N$, $A_c$ monic of degree
$g+1$), so $A_{j+1} = (-1)^g C_c(-z)$ is monic of degree $g$.
\end{proof}

\begin{lemma}[Exactness of the reduction step]
\label{lem:cusp-exact}
\uses{def:cusp-fold, lem:cusp-reduction}
\lean{ChabautyColeman.cuspNext_exact}
\leanok
$A_{j+1}(z)\,A_c(-z) - A_{j+1}(-z)\,A_c(z) = (-1)^{g+1}\, 2 z^N$, exactly.
\end{lemma}
\begin{proof}
\leanok
$(-1)^g\bigl(A_{j+1}(z)A_c(-z) - A_{j+1}(-z)A_c(z)\bigr) = C_c(-z)A_c(-z) - C_c(z)A_c(z)
= \bigl(-z^N - V_c(z^2)\bigr) - \bigl(z^N - V_c(z^2)\bigr) = -2z^N$.
\end{proof}

\begin{theorem}[The moment characterization]
\label{thm:cusp-moment}
\uses{lem:cusp-divisibility, lem:cusp-reduction, lem:cusp-exact}
\lean{ChabautyColeman.cuspMoment_step, ChabautyColeman.cuspMoment_congruence,
  ChabautyColeman.cuspMoment_base, ChabautyColeman.coeff_mul_comp_neg_X_eq_zero,
  ChabautyColeman.momentMatrix, ChabautyColeman.momentRhs,
  ChabautyColeman.momentMatrix_mulVec, ChabautyColeman.det_smul_eq_cramer_of_mulVec,
  ChabautyColeman.momentDet, ChabautyColeman.momentDet_linear,
  ChabautyColeman.moment_system_of_congruence, ChabautyColeman.momentDet_value_atom,
  ChabautyColeman.eq_of_det_momentMatrix_ne_zero}
\leanok
Let $T_j$ be the multiset $\{t_1, \dots, t_g\} \uplus \{t_1, \dots, t_j\}$, $G := g + j$,
$e_r := e_r(T_j)$, and $d_j(z) := \prod_i (t_i - z)^{\epsilon_i}$ with $\epsilon_i = 2$
for $i \le j$ and $1$ otherwise. If the avoidance values $A_i(0)$ for $i < j$ are
nonzero, then:
\begin{enumerate}
\item \textup{(master congruence)}\ \
  $A_j(z)\, d_j(-z) \equiv A_j(-z)\, d_j(z) \pmod{z^{2g+1}}$;
\item equivalently, writing $A_j = z^g + \sum_{m < g} \alpha_m z^m$: \
  $\sum_{m=0}^{g} \alpha_m\, e_{G-k+m} = 0$ for $k = 1, 3, \dots, 2g-1$
  \textup{(the \emph{moment system}; $\alpha_g := 1$)};
\item if moreover $\det M_j \neq 0$ for the system matrix
  $M_j := \bigl(e_{G-2i+1+m}\bigr)_{1 \le i \le g,\ 0 \le m \le g-1}$, then $A_j$ is the
  unique monic degree-$g$ solution, and each condition value of the fold appears as a
  pattern determinant of a once-extended alphabet. The phase-2 condition values are the
  mirror values $A_j(\mp t_j)$ of the retraced phase-1 stages
  (Lemma~\ref{lem:fold-mirror}).
\end{enumerate}
\end{theorem}
\begin{proof}
\leanok
(1) Induction. Base ($j = 0$): $d_0 = (-1)^g A_0$, so
$A_0(z)d_0(-z) = A_0(-z)d_0(z)$ holds \emph{exactly}. Composition multiplies the defect
by $(z^2 - t_{j+1}^2)$:
$A_c(z)d_c(-z) - A_c(-z)d_c(z) =
(z^2 - t_{j+1}^2)\bigl(A_j(z)d_j(-z) - A_j(-z)d_j(z)\bigr)$, where
$d_c := d_j \cdot (t_{j+1} - z)$. Reduction: by Lemma~\ref{lem:cusp-exact},
$A_{j+1}(z) \equiv A_{j+1}(-z)\,A_c(z)\,A_c(-z)^{-1} \pmod{z^{2g+1}}$, since the exact
defect $(-1)^{g+1}2z^N$ vanishes modulo $z^{2g+1}$ ($N = 2g+1$) and $A_c(-z)$ is
invertible modulo $z^{2g+1}$ precisely when $A_c(0) = -t_{j+1}A_j(0) \neq 0$ --- the
avoidance hypothesis. Hence the congruence for $A_c$ transfers to $A_{j+1}$.
(2) $\Gamma(z) := A(z)d(-z) - A(-z)d(z)$ is an odd polynomial, and for odd $k$ its
$z^k$-coefficient equals $2\sum_m \alpha_m e_{G-k+m}$ (write $d(-z) = \sum_r e_{G-r} z^r$
and collect the parity signs, which cancel); in characteristic $\neq 2$ the congruence
modulo $z^{2g+1}$ is thus exactly the system.
(3) Uniqueness is Cramer's rule; the value formulas follow from
$U_j(z^2) = (-1)^g A_j(z)A_j(-z)$ (Lemma~\ref{lem:cusp-reduction}) evaluated at
$z = t_{j+1}$, $t_j$, $0$.
\end{proof}

\begin{remark}[The Schur-polynomial interpretation of the atoms]
\label{rem:schur-atoms}
The determinant values certified in Theorem~\ref{thm:cusp-witness} by the
leading-coefficient induction admit an explicit Schur-polynomial theory. It is
\emph{deliberately left informal} --- the formalized route certifies the determinants
directly and does not use the Schur form --- and is included here as an independent
verification, checked symbolically in Magma, and as the structural explanation of the
$\pm 2^a$ integer contents observed in all numerical data.

\emph{(1) Cramer data are Schur polynomials.}
Let $\Lambda(j) := (g^{\,j+1}, g-1, g-2, \dots, 1)$. Then, for $0 \le m \le g$, the
$m$-th Cramer determinant satisfies
$\mathrm{NA}_m = (-1)^{g-m}\, s_{\Lambda(j)/(m)}(T_j)$ (a skew Schur polynomial, $(m)$
the one-row partition); in particular
$\det M_j = \pm\, s_{(g^j,\, g-1, \dots, 1)}(T_j)$ --- at $j = 0$ the classical
$s_{\delta}(t_1, \dots, t_g) = \prod_{i < l}(t_i + t_l)$ --- and by the one-variable
branching rule $\mathrm{NA}_j(-w) = \pm\, s_{\Lambda(j)}\bigl(T_j \cup \{w\}\bigr)$.
Consequently every fold condition is, up to sign and the factor $\det M_j^{2}$, a
product of two values of the single Schur polynomial $s_{\Lambda(j)}$ at multiset
alphabets: $s_{\Lambda(j)}(T_j \cup \{t\})\, \cdot\, s_{\Lambda(j)}(T_j \cup \{-t\})$
with $t = t_{j+1}$ (step) resp.\ $t = t_j$ (mirror), and $s_{\Lambda(j)}(T_j)^2$
(avoidance). The values at the all-positive alphabets are monomial-positive.
\emph{Why:} the matrix $M_j$ and its column-replaced Cramer variants are minors of the
Toeplitz array $(e_{r-c}(T_j))_{r,c}$; by the dual Jacobi--Trudi identity for skew
shapes, any such minor is $\pm$ a skew Schur polynomial, and matching the index pattern
$e_{G-2i+1+m}$ against $e_{\lambda'_i - \mu'_s - i + s}$ gives
$\lambda' = (G, G-1, \dots, G-g+1)$, i.e.\ $\lambda = \Lambda(j)$, and $\mu' = (1^m)$,
i.e.\ $\mu = (m)$; the sign $(-1)^{g-m}$ is the column transposition moving the
replaced column to its place. For $m = g$ the skew shape drops the full first row,
leaving the straight $(g^j, g-1, \dots, 1)$. The branching rule
$s_{\Lambda}(X \cup \{w\}) = \sum_{r \ge 0} s_{\Lambda/(r)}(X)\, w^r$ (one extra
variable picks up a horizontal strip, necessarily a single row) gives the stated form of
$\mathrm{NA}_j(-w)$; monomial-positivity of (skew) Schur polynomials is
Lindstr\"om--Gessel--Viennot.

\emph{(2) Nested-extreme factorization.}
Let $\Lambda$ be a partition, $\ell(\Lambda) \le n$. Partition the $n$ slots of
$s_\Lambda(x_1, \dots, x_n)$ into $k$ groups of sizes $m_1, \dots, m_k$ carrying
distinct values $v_1 > \dots > v_k$ (a priority order), and let $S(v_1, \dots, v_k)$ be
the resulting polynomial. Define $D_r := \deg_{v_r}$ and the iterated leading
coefficients along the priority order. Then, with $M_r := m_1 + \dots + m_r$ and the
\emph{slices} $\sigma_r := (\Lambda_{M_{r-1}+1}, \dots, \Lambda_{M_r})$:
$D_r = |\sigma_r|$, and the final iterated leading coefficient equals
$\prod_r N_r$ with $N_r := s_{\sigma_r}(1^{m_r}) \ge 1$, the number of semistandard
Young tableaux of the straight shape $\sigma_r$ with $m_r$ letters. Equivalently: the
monomial $v_1^{D_1} \cdots v_k^{D_k}$ is the lexicographically largest monomial of $S$
and has coefficient $\prod_r N_r$.
\emph{Why:} by symmetry assign group $r$ the entry range $\{M_{r-1}+1, \dots, M_r\}$;
the $v_1$-degree of a tableau's monomial is its number of entries $\le m_1$. In a
semistandard tableau the entry in row $i$ is $\ge i$, so entries $\le m_1$ lie in rows
$\le m_1$ and form a Young subdiagram $\mu \subseteq \Lambda$ (left-closed by
row-weakness, up-closed by column-strictness); thus
$|\mu| \le \Lambda_1 + \dots + \Lambda_{m_1}$ with equality iff $\mu$ is the full top
$m_1$ rows, and the bound is attained. A tableau attaining it is exactly a pair: an
arbitrary tableau of the straight shape $\sigma_1$ with $m_1$ letters ($N_1$ choices),
and an arbitrary filling of the lower rows with the larger entries; the lower part, with
entries shifted down by $m_1$, is an arbitrary tableau of the straight shape
$(\Lambda_{m_1+1}, \dots)$ on the remaining letters. Induction on $k$ finishes;
$N_r \ge 1$ since $\ell(\sigma_r) \le m_r$.

\emph{(3) Slice values.}
By Weyl's dimension formula, $s_{(a)}(1) = 1$, $s_{(a,b)}(1,1) = a - b + 1$ and
\[ s_{(a,b,c)}(1,1,1) = \tfrac12\,(a-b+1)(b-c+1)(a-c+2)\,. \]
Since consecutive rows of $\Lambda(j)$ differ by $0$ or $1$: every $2$-slice contributes $1$ or $2$; a
$3$-slice contributes $1$ on $(a,a,a)$, $8$ on $(a, a{-}1, a{-}2)$, and $3$ on the
boundary patterns $(a,a,a{-}1)$, $(a,a{-}1,a{-}1)$ --- the source of the
characteristic-$3$ fingerprint observed in the numerical data.

\emph{(4) Certificates for the positive atoms.}
Each of the values $s_{\Lambda(j)}(T_j)$, $s_{\Lambda(j)}(T_j \cup \{t_{j+1}\})$,
$s_{\Lambda(j)}(T_j \cup \{t_j\})$ ($j \ge 1$, $g \ge 2$) contains an explicit monomial
with coefficient a power of $2$; consequently each is nonzero over every field of
characteristic $\neq 2$, and has $2$-power integer content. Specifically: for the first
two alphabets (doubles and singles only) \emph{every} priority order works, with
$2$-exponent the number of difference-$1$ two-row slices; for the tripled alphabet,
place the triple first if $j \ge 2$ (slice $(g,g,g) \mapsto 1$) and after one single if
$j = 1$ (slice $(g, g{-}1, g{-}2) \mapsto 8$). This is (2) combined with (3).

\emph{(5) Signed group values.}
(a) $s_\mu(t, -t) = 0$ unless $\ell(\mu) \le 2$ and $\mu_1 \equiv \mu_2 \bmod 2$, in
which case it equals $(-1)^{\mu_2} t^{|\mu|}$; in particular
$s_{(g,g)}(t,-t) = (-1)^g t^{2g}$.
(b) $s_\mu(t, t, -t) = t^{|\mu|} s_\mu(1,1,-1)$, and
$s_{(a,a,a)}(1,1,-1) = s_{(a,a,a-1)}(1,1,-1) = (-1)^a$.
(For (a), $s_{(\mu_1,\mu_2)}(x,y) = (xy)^{\mu_2} h_{\mu_1-\mu_2}(x,y)$ and
$h_d(t,-t) = t^d$ for even $d$, $0$ for odd $d$. For (b),
$s_{(a,a,a)} = (xyz)^a \mapsto (-1)^a$;
$s_{(a,a,a-1)} = (xyz)^{a-1} e_2 \mapsto (-1)^{a-1}\cdot(1 - 1 - 1) = (-1)^a$. Note the
contrast with the \emph{unsigned} $s_{(a,a,a-1)}(1,1,1) = 3$: the signed specialization
rescues exactly the boundary pattern.)

\emph{(6) Certificates for the negated atoms.}
For $j \ge 1$, $g \ge 2$ and $\Lambda = \Lambda(j)$:
in $s_{\Lambda}(T_j \cup \{-t_{j+1}\})$ the $t_{j+1}$-degree is exactly $2g$, with
leading coefficient $(-1)^g\, s_{\Lambda''}\bigl(T_j \smallsetminus \{t_{j+1}\}\bigr)$,
$\Lambda''$ the shape with the top two rows removed; in
$s_{\Lambda}(T_j \cup \{-t_j\})$ the $t_j$-degree is exactly $3g$ ($j \ge 2$) resp.\
$3g - 1$ ($j = 1$), with leading coefficient $(-1)^g\, s_{\Lambda'''}(\text{rest})$,
$\Lambda'''$ the shape with the top three rows removed. Continuing with (2) on the
remaining alphabet, each negated atom contains an explicit monomial with coefficient
$\pm 2^a$, hence is nonzero over every field of characteristic $\neq 2$ with $2$-power
content. \emph{Why:} give the signed group --- the pair $\{t, -t\}$ (step case) resp.\
the trio $\{t, t, -t\}$ (mirror case) --- top priority and expand by the coproduct
$s_\Lambda(Y \uplus X) = \sum_\mu s_\mu(Y)\, s_{\Lambda/\mu}(X)$, so the $t$-degree
comes only from $s_\mu(Y)$. In the pair case
$|\mu| \le \Lambda_1 + \Lambda_2 = 2g$ with equality only for $\mu = (g,g)$, which
passes the parity condition of (5)(a) with value $(-1)^g$; in the trio case
$|\mu| \le \Lambda_1 + \Lambda_2 + \Lambda_3$ with equality only for the full top three
rows, with value $(-1)^g$ by (5)(b). In each case $\mu$ consists of full equal-length
top rows, so $\Lambda/\mu$ is the straight shape with those rows removed, its
consecutive-difference pattern is again $0/1$, and no larger $|\mu|$ contributes.

The statements of (1)--(6) were verified symbolically for small $g$ (scripts
\texttt{schurid.m}, \texttt{p2shapes.m}, \texttt{branch.m}, \texttt{p34atoms.m},
\texttt{c3check.m} in \texttt{prototype/}). The formalized counterparts of (4) and (6)
are the leading-coefficient cascades cited with Theorem~\ref{thm:cusp-witness}; the
degree and leading-coefficient statements of (6) are exactly the coefficient identities
of the peeled node there.
\end{remark}

\begin{theorem}[The cuspidal witness, all $g$]
\label{thm:cusp-witness}
\lean{ChabautyColeman.cusp_alphabet_comp_neg,
  ChabautyColeman.momentDet_mul, ChabautyColeman.peelPattern,
  ChabautyColeman.peelStage_single, ChabautyColeman.peelStage_double,
  ChabautyColeman.peelStage_id, ChabautyColeman.peelStage_triple_id,
  ChabautyColeman.peelStage_lead,
  ChabautyColeman.coeff_momentDet_of_stage, ChabautyColeman.degree_prod_coeff_lt,
  ChabautyColeman.coeff_momentDet_single, ChabautyColeman.coeff_momentDet_double,
  ChabautyColeman.coeff_momentDet_id, ChabautyColeman.coeff_momentDet_triple_id,
  ChabautyColeman.coeff_momentDet_lead, ChabautyColeman.peelHom,
  ChabautyColeman.momentDet_ne_zero_of_map,
  ChabautyColeman.momentDet_singles_ne_zero, ChabautyColeman.momentDet_doubles_ne_zero,
  ChabautyColeman.momentDet_X_peelPattern_one,
  ChabautyColeman.eval_ne_zero_of_momentDet_ne_zero,
  ChabautyColeman.momentDet_generic_doubles_ne_zero,
  ChabautyColeman.momentDet_generic_doubles_mul_X_ne_zero,
  ChabautyColeman.momentDet_generic_pair_ne_zero,
  ChabautyColeman.momentDet_generic_trio_ne_zero,
  ChabautyColeman.momentDet_generic_tripled_ne_zero,
  ChabautyColeman.momentDet_generic_trio_id_ne_zero,
  ChabautyColeman.momentDet_generic_tripled_one_ne_zero,
  ChabautyColeman.cuspGenH, ChabautyColeman.cuspChain_cond_of_atoms,
  ChabautyColeman.isCuspState_cuspChain_of_atoms,
  ChabautyColeman.cuspChain_congruence_of_atoms,
  ChabautyColeman.cuspChain_eval_ne_zero_of_atoms,
  ChabautyColeman.cuspStep_cond_base, ChabautyColeman.exists_cuspStep_H,
  ChabautyColeman.genNode, ChabautyColeman.genCuspChain,
  ChabautyColeman.genNode_trio_atom_ne_zero,
  ChabautyColeman.genNode_tripled_atom_ne_zero,
  ChabautyColeman.exists_genCuspChain_witness, ChabautyColeman.genFoldH,
  ChabautyColeman.genCuspFold, ChabautyColeman.exists_genCuspFold_witness,
  ChabautyColeman.genCuspFold_eval_sq_ne_zero,
  ChabautyColeman.isScaledState_foldBase_gen, ChabautyColeman.genFoldH_eq_foldGenH,
  ChabautyColeman.foldCondNum_cusp_ne_zero, ChabautyColeman.cuspNum_ne_zero}
\leanok
\uses{thm:cusp-moment, lem:fold-mirror}
For every $g \ge 2$ and every field $k$ of characteristic $\neq 2$: at the generic
nodes $t_1, \dots, t_g$ the full $2g$-step cuspidal fold is defined, with every
composition, disjointness and avoidance condition of both phases nonzero in
$k(t_1, \dots, t_g)$, and the moment determinants $\det M_j$ are nonzero for all
$j \le g$. Consequently the cleared condition numerators $P_q$ ($q < 2g$) of the
numerator fold at the cuspidal inputs are nonzero polynomials over $k$ --- the
nonvanishing hypothesis that Lemma~\ref{lem:cusp-descent} consumes.
\end{theorem}
\begin{proof}
\leanok
The values to be certified --- the determinants $\det M_j$, the avoidance values
$\mathrm{NA}_j(0)$, the step values $\mathrm{NA}_j(\mp t)$ and the mirror values
$\mathrm{NA}_j(\mp t_j)$ at the last-added node --- are nonzero at the generic nodes by
an exact leading-coefficient induction, one node at a time. Read the pattern determinant
in the polynomial ring in the peeled node $u$ over the remaining generic nodes. For each
of the five factor shapes occurring in the alphabets (a linear factor, its square, a
signed pair, a trio, and a tripled factor), the identity of the corresponding peel stage
states that the top $u$-coefficient of the determinant is a fixed scalar --- $1$, $2$,
$\pm 1$ or $8$ --- times the determinant of the remaining alphabet at the next pattern,
and no other term of the expansion reaches that degree. Nonvanishing therefore follows
by induction on the number of nodes over the integral domain of generic nodes; the
extension atoms peel their distinguished node first. No valuation-theoretic input
enters. The chain induction built on Theorem~\ref{thm:cusp-moment} then discharges
every phase-1 condition from these determinant values, and the mirror
(Lemma~\ref{lem:fold-mirror}) extends the induction through phase 2 along the retraced
phase-1 alphabets. Identifying the resulting scalars with the generic scalars of the
numerator fold and clearing denominators converts each condition into its numerator
$P_q$.

The Schur-side certificates of Remark~\ref{rem:schur-atoms} give an independent
verification; the numerical computation of every condition for $g \le 6$ from the moment determinants
(script \texttt{detcond.m}) agrees with the exact fold wherever both are available.
\end{proof}

\begin{lemma}[Descent from the cusp]
\label{lem:cusp-descent}
\uses{lem:fold-universal, thm:cusp-witness}
\lean{ChabautyColeman.descRate, ChabautyColeman.digits_eq_of_sum_eq,
  MvPolynomial.IsWeightedHomogeneous.aeval,
  ChabautyColeman.orderTop_aeval_descRate_le, ChabautyColeman.le_orderTop_mul_sub_one,
  ChabautyColeman.le_orderTop_prod_sub_one, ChabautyColeman.aeval_family_ne_zero,
  ChabautyColeman.exists_infinityPoint}
\leanok
Fix a field $K$ of characteristic $\neq 2$. Let $\mathcal{D} \ge 1$ bound the
$t$-degrees of the finitely many cuspidal numerators $P_q$ of
Lemma~\ref{lem:fold-universal}, choose a base $B > \mathcal{D}$, set
$M := \mathcal{D}B^{\,g+1}$, and use the \emph{common rates}
$\tilde n_i := M + B^{\,i+1}$ for every fold quantity $q$. If $P_q \neq 0$, then the
value of the universal numerator $\mathcal{N}_q$ at the family of points escaping to
infinity with rates $\tilde n_i$ is a \emph{nonzero} element of $K(\!(s)\!)$: its
valuation is that of the valuation-minimal monomial of
$P_q(t)\rvert_{t_i = s^{-\tilde n_i}}$, which lies in the window
$[\,-dM - \mathcal{D}B^{\,g},\, -dM\,]$ for $d = \operatorname{wt}(\mathcal{N}_q)$.
\end{lemma}
\begin{proof}
\leanok
\emph{(a) Separation and the window.} $P_q$ is homogeneous of some degree
$d \le \mathcal{D}$, so a monomial $t^c$ of $P_q$ specializes to valuation
$-\sum_i c_i \tilde n_i = -dM - \sum_i c_i B^{\,i+1}$: the $M$-part is common to
\emph{all} monomials (homogeneity), and since $0 \le c_i \le d \le \mathcal{D} < B$ the
digits $c_i$ make $c \mapsto \sum_i c_i B^{\,i+1}$ injective. Hence distinct monomials
land on distinct valuations, all inside the window
$[\,-dM - \mathcal{D}B^{\,g},\, -dM\,]$, and
$v^\ast := v\bigl(P_q(s^{-\tilde n})\bigr)$ is attained by a unique monomial; in
particular $P_q(s^{-\tilde n}) \neq 0$.

\emph{(b) The correction budget.} At the family, $y_i = t_i^N (1 + \eta_i)$ with
$v(\eta_i) \ge 2\tilde n_i \ge 2M$ (a Hensel square root), and $x_i = t_i^2$ exactly.
Evaluate $\mathcal{N}_q$ monomial by monomial. The deviation of an $a$-free monomial
$\mu$ from its cuspidal contribution has valuation
$\ge v(\mu(\text{cusp})) + 2M \ge (-dM - \mathcal{D}B^{\,g}) + 2M > -dM \ge v^\ast$,
using that \emph{every} $a$-free monomial of the isobaric $\mathcal{N}_q$ --- whether
or not it survives the collection into $P_q$ --- has $t$-degree exactly $d$ and hence
lies in the window of (a). A monomial containing some $a_r$ has $(x,y)$-weight at most
$d - 2$, so its value has valuation
$\ge -(d-2)(M + B^g) \ge -dM + 2M - \mathcal{D}B^{\,g} > -dM \ge v^\ast$. Hence
$v\bigl(\mathcal{N}_q(\text{family}) - P_q(s^{-\tilde n})\bigr) > v^\ast$, so
$\mathcal{N}_q(\text{family})$ has valuation exactly $v^\ast$; in particular it is
nonzero.
\end{proof}

\begin{theorem}[The fold witness for genus $\ge 2$]
\label{thm:fold-witness}
\lean{ChabautyColeman.exists_foldWitness, ChabautyColeman.curveCondNum,
  ChabautyColeman.curveSpec, ChabautyColeman.not_curveVanishing_curveCondNum,
  ChabautyColeman.eval_curveCondNum, ChabautyColeman.foldXs, ChabautyColeman.foldYs,
  ChabautyColeman.isScaledState_foldBase, ChabautyColeman.foldCond_ne_zero}
\leanok
\uses{lem:fold-universal, thm:cusp-witness, lem:cusp-descent}
Let $g \ge 2$ and let $K$ be algebraically closed of characteristic $\neq 2$, $f$ monic
of degree $N = 2g+1$. Every universal fold numerator $\mathcal{N}_q$ is a nonzero
polynomial function on $g$-tuples of affine curve points; consequently there is a
\emph{single} tuple $P_1, \dots, P_g$ over $K$, with pairwise distinct $x$-coordinates,
at which all $\mathcal{N}_q$ are simultaneously nonzero --- so the point-by-point fold
$\tau + \tau - \tau$ at this tuple is defined throughout and every disjointness and
doubling condition of both phases is nonzero, for every genus $\ge 2$. (Genus $1$,
where the cuspidal specialization degenerates in characteristic $3$, is
Lemma~\ref{lem:fold-witness-genus-one}.)
\end{theorem}
\begin{proof}
\leanok
Each $\mathcal{N}_q$ is nonzero on the variety of point tuples: the family of points
escaping to infinity with the common rates of Lemma~\ref{lem:cusp-descent} is a
$K(\!(s)\!)$-point at which $\mathcal{N}_q$ takes a nonzero value --- a pure polynomial
evaluation, requiring no other condition to hold, hence no circularity. A common
witness tuple then exists because finitely many nonzero polynomial conditions on $C^g$
have a common witness (a product of the conditions is refuted by one witness per
factor); the pairwise distinctness of the $x$-coordinates joins as $g(g-1)$ further
conditions $x_i - x_j$, each refuted at a two-point tuple. At the common tuple the fold
quantities are the ratios of Lemma~\ref{lem:fold-universal} with nonvanishing
numerators and denominators, inductively along the fold. Phase 2 needs no separate
argument: its conditions belong to the same uniform family of numerators, their
cuspidal nonvanishing having been established through the mirror inside
Theorem~\ref{thm:cusp-witness}. Further polynomial avoidance conditions --- the
Weierstrass avoidance $\operatorname{res}(U_j, f) \neq 0$ of the intermediate states,
or $y_i \neq 0$ --- join the same argument by enlarging the finite set of conditions,
their cuspidal nonvanishing being recorded there as well (on the cusp
$\operatorname{res}(U_j, X^N) = \pm U_j(0)^N$ and $U_j(0) = \pm A_j(0)^2 \neq 0$).
\end{proof}

\begin{lemma}[The genus-one fold witness]
\label{lem:fold-witness-genus-one}
\lean{ChabautyColeman.exists_genusOnePoint,
  ChabautyColeman.GenusOne.tangent_factorization,
  ChabautyColeman.GenusOne.foldChain_one, ChabautyColeman.GenusOne.foldChain_two,
  ChabautyColeman.GenusOne.foldGenH_zero_eq, ChabautyColeman.GenusOne.foldGenH_one_eq,
  ChabautyColeman.GenusOne.psi_eval_eq, ChabautyColeman.GenusOne.avoid_eval_eq,
  ChabautyColeman.GenusOne.psi_ne_zero, ChabautyColeman.GenusOne.avoid_ne_zero}
\leanok
\uses{lem:fold-universal}
Let $g = 1$ and let $K$ be algebraically closed of characteristic $\neq 2$, $f$ monic
cubic and \emph{squarefree}. Some affine curve point $(x_0, y_0)$ has $y_0 \neq 0$,
nonzero step-condition numerators at both fold steps, and nonzero stage resultants at
all three fold stages --- the genus-one instance of Theorem~\ref{thm:fold-witness}.

The cuspidal argument does not cover this case: on the cusp the interior condition
value at $g = 1$ is $3t^2/4$, identically zero in characteristic $3$ (the additive
group of the cusp is all $3$-torsion there, so $x(2P) = x(P)$ throughout), and for
$f = X^3$ the conditions genuinely vanish at every curve point --- squarefreeness is
sharp.
\end{lemma}
\begin{proof}
\leanok
A direct tangent--chord computation of the fold
(Lemma~\ref{lem:fold-universal} supplies the transfer to the numerator conditions).
With $\lambda = f'(x_0)/(2y_0)$, $x_2 = \lambda^2 - a - 2x_0$ and
$y_2 = -y_0 - \lambda(x_2 - x_0)$, the single cubic factorization
\[ (y_0 + \lambda(X - x_0))^2 - f \;=\; -(X - x_0)^2 (X - x_2) \]
governs both steps: the chord through $2P$ and $-P$ is the \emph{negated} tangent line,
so the second step needs no new identity, and the chain visits $(X - x_2, y_2)$ and
returns to $(X - x_0, y_0)$ --- $\tau + \tau - \tau = \tau$ on the nose. The condition
values are, up to powers of $2y_0$, the values at $x_0$ of $\psi_3 = 2ff'' - f'^2$ (the
$3$-division polynomial: $\psi_3(x_0) = 4y_0^2\,(x_0 - x_2)$) and of
$f'\psi_3 - 8f^2$ (the numerator of $y(2P)$: $8y_0^3\,y_2$); the stage resultants are
the $y$-values $y_0, y_2, y_0$ of the states, again up to denominators. Both
polynomials are nonzero for squarefree $f$ in \emph{every} characteristic, by
separability: if either vanished, $f \mid f'^2$, hence $f \mid f'$ by coprimality,
hence $f' = 0$ by degrees --- contradicting the coprimality of $f$ and $f' = 0$ for the
nonunit $f$. The witness point takes $x_0$ outside the finitely many roots of
$f \cdot \psi_3 \cdot (f'\psi_3 - 8f^2)$ and $y_0 = \sqrt{f(x_0)}$.
\end{proof}

\begin{lemma}[The auxiliary class $\tau$]
\label{lem:tau-exists}
\lean{ChabautyColeman.TauData, ChabautyColeman.TauData.exists_of_hasGoodCenter,
  ChabautyColeman.exists_interpolating_datum,
  ChabautyColeman.exists_ftauUY_root, ChabautyColeman.exists_isFoldCerts_lift,
  ChabautyColeman.exists_isFoldCerts_ftau, ChabautyColeman.isFtauFoldAvoid_of_center,
  ChabautyColeman.isFoldAvoid_of_map, ChabautyColeman.exists_bezout_lift_map,
  ChabautyColeman.TauData.exists_genus_two}
\leanok
\uses{lem:tau-generic, lem:many-points, lem:disk-param, lem:mumford-exists}
There exist $q$ and a class $\tau \in \J{\Qq}$ whose reduced datum $(u_\tau, v_\tau)$
lifts the witness datum of Lemma~\ref{lem:tau-generic} and satisfies
Definition~\ref{def:goodpos}, together with all the data the constructions below
require: an integral Mumford datum with its cofactor identity and a B\'ezout identity
$c_1 u_\tau + c_2 (2 v_\tau) = 1$ over $\Oq$; the $g$ points splitting it,
$u_\tau = \prod_i (X - a_i)$ with pairwise unit differences $a_i - a_j$ and
$v_\tau(a_i) = b_i$; the \emph{moving points} of the fold --- the Hensel roots of the
perturbed $u$-polynomial over $\Oq[[X;Y]]$, with constant coefficients the $a_i$ ---
and valid B\'ezout data for the fold computation of $F_\tau$ over $\Oq[[X;Y]]$,
avoiding the Weierstrass points at every stage. The residue-level genericity of
Lemma~\ref{lem:tau-generic} travels with the datum: the lifts have prescribed residues,
and the avoidance conditions descend along the residue map. The avoidance of the
prescribed finite set $T$ of residue $x$-values --- $\bar u_\tau(x) \neq 0$ for
$x \in T$ --- also comes with the datum; the per-disk construction of
Lemma~\ref{lem:global-generic} consumes it.
\end{lemma}
\begin{proof}
\leanok
The witness tuple of Lemma~\ref{lem:tau-generic} lifts point by point, by Hensel's
lemma, to disk centres with prescribed residues (Lemma~\ref{lem:disk-param}); $u_\tau$
is the product of the lifted linear factors, and $v_\tau$ is interpolated through the
lifted points by the Chinese remainder theorem --- the lifted $x$-coordinates differ by
units, so the linear factors are pairwise comaximal --- with $\bar v_\tau$ recovered
exactly by uniqueness of interpolation in degree $< g$. The B\'ezout identity over
$\Oq$ lifts by Nakayama's lemma from the stage-$0$ avoidance of the witness. The moving
points exist by Hensel's lemma along the ideal generated by the variables: the
$\tau$-roots are simple roots of $u_\tau$, their derivative values being unit products
of the pairwise differences. The B\'ezout data for the fold lift step by step, again by
Nakayama with prescribed residues against the already-lifted states, and the
Weierstrass avoidance descends from the residue configuration, the residue map
reflecting coprimality of monic polynomials.
\end{proof}

\begin{definition}[Cantor addition as a rational map]
\label{def:cantor-map}
\lean{ChabautyColeman.cantorV, ChabautyColeman.cantorV_sq_sub,
  ChabautyColeman.y_sub_algebraMap_cantorV, ChabautyColeman.cantorCompose,
  ChabautyColeman.cantorReduce}
\leanok
\uses{lem:mumford-exists}
The pair of tuples of rational functions
$\mathrm{Add} = (\mathrm{Add}_u, \mathrm{Add}_v)$ in the coordinates of two Mumford
data of degree $g$, obtained by transcribing Cantor's composition and reduction
symbolically; its domain of definition contains the locus where the finitely many
leading coefficients and resultants appearing in the algorithm are invertible. The
inverse is $\ominus(u, v) := (u, -v)$ (Lemma~\ref{lem:aj-inverse}). The composition is
the coprime branch of Cantor's algorithm with the three-term B\'ezout relation
$s_1 u_1 + s_2 u_2 + s_3(v_1 + v_2) = 1$: unlike the extended Euclidean algorithm,
which is not a rational map near the diagonal, this branch uniformly covers
composition, the near-diagonal case, and doubling.
\end{definition}

\begin{lemma}[Correctness of the rational map]
\label{lem:cantor-correct}
\lean{ChabautyColeman.mumfordIdeal_mul_of_bezout,
  ChabautyColeman.mk0_mumfordIdeal_mul_of_bezout,
  ChabautyColeman.mk0_mumfordIdeal_eq_mk0_cofactor_datum,
  ChabautyColeman.exists_bezout_lift}
\leanok
\uses{def:cantor-map, lem:mumford-exists}
At every point of its domain over any field, $\mathrm{Add}$ evaluates to the reduced
datum of the product class. More precisely, in the certificate-passing form used
throughout: given a B\'ezout identity $s_1 u_1 + s_2 u_2 + s_3(v_1+v_2) = 1$, the
product of the two Mumford ideals \emph{equals} the Mumford ideal of the composed datum
$(u_1 u_2,\ V)$ with $V = s_1 u_1 v_2 + s_2 u_2 v_1 + s_3(v_1 v_2 + f)$, so the
composition computes the product class; and given an inverse of the leading coefficient
of the cofactor $w$ of $v^2 - f = uw$, the normalized reduced datum
$(c\,w,\ (-v) \bmod c\,w)$ is monic, is again a Mumford datum, and represents the same
class (Lemma~\ref{lem:mumford-reduce}). Over a local ring, a B\'ezout identity for the
residues lifts by Nakayama's lemma, since the Mumford $u$-polynomials are monic
(monicity is essential: $1 - \pi X$ over a discrete valuation ring has invertible
residue but is not invertible).
\end{lemma}
\begin{proof}
\leanok
The composed ideal identity is proved directly from the certificate:
$V \equiv v_i \bmod u_i$ and $V^2 \equiv f \bmod u_1 u_2$ with explicit cofactors, and
$y - V$ is an explicit combination of the four product generators; the reduction step
reduces to Lemma~\ref{lem:mumford-reduce} after normalizing by the leading-coefficient
inverse.
\end{proof}

\begin{definition}[The chart at $\tau$]
\label{def:chart}
\lean{ChabautyColeman.chartU, ChabautyColeman.chartV, ChabautyColeman.exists_chartV}
\leanok
\uses{lem:tau-exists, thm:hensel}
Coordinates near $\tau$: the $g$ coefficients $X = (X_1,\dots,X_g)$ of $u - u_\tau$.
There is $V(X) \in \Oq[[X]][x]$ of $x$-degree $< g$ with $V(0) = v_\tau$ and
$V(X)^2 \equiv f \bmod (u_\tau + X)$; the chart point of $X$ is the Mumford datum
$(u_\tau + X,\ V(X))$, i.e.\ the ideal $(u_\tau + X,\ y - V(X))$. \emph{Formalization.}
$V$ is constructed as a single quadratic Hensel root (Theorem~\ref{thm:hensel}) of
$T^2 = f$ in the quotient by $(u_\tau + X)$, a finite free algebra over the power
series ring, which is Henselian for the $(X)$-adic ideal; no topology on $\Oq$ is
needed. The inputs are the cofactor identity $v_\tau^2 - f = u_\tau w$ and the B\'ezout
identity $c_1 u_\tau + c_2 (2 v_\tau) = 1$ carried by the datum of
Lemma~\ref{lem:tau-exists}.
\end{definition}

\begin{definition}[The translated group law $F_\tau$]
\label{def:Ftau}
\lean{ChabautyColeman.FtauFold, ChabautyColeman.ftauFoldPipe,
  ChabautyColeman.FoldCertData, ChabautyColeman.IsFoldCerts,
  ChabautyColeman.IsFtauFoldCerts, ChabautyColeman.ftauFoldPipe_fst_eq,
  ChabautyColeman.ftauFoldXs, ChabautyColeman.ftauFoldYs,
  ChabautyColeman.ftauUY_eq_prod}
\leanok
\uses{def:chart, def:cantor-map, lem:tau-exists}
$F_\tau(X,Y)$ is defined as the chart coordinates of
$\mathrm{ch}(X) + \mathrm{ch}(Y) - \tau$, expanded as a formal power series at $(0,0)$:
in the $\tau$-translated identification $c(X) := \mathrm{ch}(X) - \tau$ this is the
group law, satisfying $\mathrm{ch}(F_\tau(X,Y)) = \mathrm{ch}(X) + \mathrm{ch}(Y) -
\tau$. \emph{Formalization.} The computation is realized as the point-by-point fold of
Definition~\ref{def:centre-fold} over $\Oq[[X;Y]]$: starting from the $X$-side chart
datum, fold in, one point per step, first the $g$ moving points of the $Y$-side chart
--- the Hensel roots of the $Y$-side chart polynomial, which factor it completely ---
and then the $g$ involuted $\tau$-points $(a_i, -b_i)$; each step is one Cantor
composition at a B\'ezout triple followed by the unconditional reduction with leading
coefficient $-1$. The B\'ezout data are a parameter of the definition, one triple per
step; given their validity, the output is a Mumford datum whose $u$-part is exactly the
chart polynomial at $F_\tau(X,Y)$. Existence of valid data over $\Oq[[X;Y]]$ is part of
Lemma~\ref{lem:tau-exists}.
\end{definition}

\begin{lemma}[Integrality, unit laws, commutativity]
\label{lem:Ftau-basic}
\lean{ChabautyColeman.setYZero_FtauFold, ChabautyColeman.setXZero_FtauFold,
  ChabautyColeman.rename_swap_FtauFold, ChabautyColeman.constantCoeff_FtauFold,
  ChabautyColeman.fold_pipeline_eq_self, ChabautyColeman.mk0_foldPipe,
  ChabautyColeman.prod_mk0_foldXs_foldYs, ChabautyColeman.mk0_mumfordIdeal_mul_neg}
\leanok
\uses{def:Ftau, lem:cantor-correct}
$F_\tau \in \Oq[[X;Y]]^g$ (the denominators of Definition~\ref{def:cantor-map} are
units at the centre by Definition~\ref{def:goodpos}); $F_\tau(X,0) = X$,
$F_\tau(0,Y) = Y$, $F_\tau(X,Y) = F_\tau(Y,X)$; and for all
$\mathbf x, \mathbf y \in (\mathfrak m')^g$, $\mathfrak m'$ the maximal ideal of a
complete discrete valuation ring $\Oq'$ over $\Oq$,
$\mathrm{ch}(F_\tau(\mathbf x,\mathbf y))$ represents the class
$\mathrm{ch}(\mathbf x) + \mathrm{ch}(\mathbf y) - \tau$: $F_\tau$ computes the sum in
the $\tau$-translated identification.
\end{lemma}
\begin{proof}
\leanok
Integrality holds by construction, and the vanishing constant term expresses that the
value of the fold at the centre is the $\tau$-datum itself. For the unit law
$F_\tau(X,0) = X$: at $Y = 0$ the moving points degenerate onto the $\tau$-points, so
the fold computes $[\mathrm{ch}(X)] + [\tau] - [\tau] = [\mathrm{ch}(X)]$ --- the
mirrored point block cancels and reduced representatives are unique
(Lemma~\ref{lem:cantor-correct}, Lemma~\ref{lem:mumford-exists}); the identity of
power series follows by evaluating at all points and interpolating
(Theorem~\ref{thm:interpolation}; interpolation over $\Oq$-points suffices).
Commutativity is proved the same way: the swapped computations have the same class,
hence the same reduced output. The last clause is the pointwise class computation of
the fold, transported to evaluation points.
\end{proof}

\begin{theorem}[Associativity]
\label{thm:Ftau-assoc}
\lean{ChabautyColeman.FtauFold_assoc, ChabautyColeman.eval_FtauFold_eq_of_mk0_eq,
  ChabautyColeman.fold_pipeline_eq_of_mk0_eq, ChabautyColeman.mk0_chart_eval_fold,
  ChabautyColeman.chartClass, ChabautyColeman.eval_ftauFoldPipe_eq,
  ChabautyColeman.ftauFoldPipe_constantCoeff_eq, ChabautyColeman.constantCoeff_FtauFold,
  ChabautyColeman.eq_of_dvd_sq_sub_of_bezout}
\leanok
\uses{lem:Ftau-basic, thm:interpolation, lem:cantor-correct}
$F_\tau(F_\tau(X,Y),Z) = F_\tau(X,F_\tau(Y,Z))$; hence $F_\tau$ is a $g$-dimensional
commutative formal group law over $\Oq$ in the sense of Definition~\ref{def:fgl}.
\end{theorem}
\begin{proof}
\leanok
Both sides are power series over $\Oq$ agreeing at every point of the polydisk, because
at points both compute the associative class-group sum
$\mathrm{ch}(x)+\mathrm{ch}(y)+\mathrm{ch}(z)-2\tau$; conclude by
Theorem~\ref{thm:interpolation}. The pointwise statement rests on: the chart at the
evaluated coordinates $F_\tau(\mathbf x, \mathbf y)$ represents
$\mathrm{ch}(\mathbf x) + \mathrm{ch}(\mathbf y) - \tau$ (the evaluated fold output is
the chart datum at the evaluated coordinates --- the $u$-parts agree by the shape of
the result, the $v$-parts by uniqueness of square roots with a B\'ezout certificate);
and computations with class-equal inputs have equal reduced outputs, by uniqueness of
reduced representatives.
\end{proof}

\begin{lemma}[Reduction commutes with reduced representatives at good position]
\label{lem:reduce-commutes}
\lean{ChabautyColeman.red_reducedRep_eq, ChabautyColeman.red_reducedRep_points,
  ChabautyColeman.prod_pointIdeal_eq_reduced,
  ChabautyColeman.eq_of_mumfordIdeal_eq_prod_pointIdeal, ChabautyColeman.redres_sq_curve}
\leanok
\uses{thm:redJ-hom, lem:mumford-exists, def:goodpos}
Let $c \in \J{\Qq}$ be such that $\redJ(c)$ has a reduced datum
$(\bar u_0, \bar v_0)$ over $\Fq$ with $\bar u_0$ squarefree of full degree $g$. Then
the reduced datum $(u, v)$ of $c$ has coefficients in $\Oq$, $\deg u = g$, and
$(\bar u, \bar v) = (\bar u_0, \bar v_0)$.
\end{lemma}
\begin{proof}
\leanok
Compute $\redJ(c)$ from the points of the reduced representative of $c$
(Definition~\ref{def:redJ}): the split Mumford representative over the splitting field
$L$ of $u$ factors into point ideals, and the factorization determines the datum by the
norm identity $\mathfrak a\,\iota(\mathfrak a) = (u)$. Reduce it to the residue curve
and match against $(\bar u_0, \bar v_0)$ by semi-reduction and uniqueness of reduced
representatives (Lemma~\ref{lem:mumford-exists} over $\Fq$): since $\bar u_0$ has full
degree $g$ and is squarefree, no cancellation and no contribution from the disk at
infinity can occur, forcing every point of the representative to be integral with
multiplicity one and distinct reductions. Integrality of $(u, v)$ follows: $u$ has
integral roots, and $v$ is the interpolant through the integral points, the
$x$-differences being units.
\end{proof}

\begin{theorem}[Chart parametrization of the kernel of reduction]
\label{thm:chart-mem}
\lean{ChabautyColeman.J1_equiv_polydisk, ChabautyColeman.polydiskToKer,
  ChabautyColeman.polydiskToKer_bijective, ChabautyColeman.polydiskToKer_eval_FtauFold,
  ChabautyColeman.exists_chartClass_eq, ChabautyColeman.chartClass_injective,
  ChabautyColeman.redJ_chartClass, ChabautyColeman.redJ_tauClass,
  ChabautyColeman.redJ_mk0_mumford, ChabautyColeman.mumfordIdeal_eq_prod_pointIdeal}
\leanok
\uses{thm:redJ-hom, lem:reduce-commutes, thm:hensel, def:chart, lem:Ftau-basic}
The map $c \mapsto{}$chart coordinates of $c + \tau$ is a group isomorphism
$J^1(\Qq) \xrightarrow{\ \sim\ } \big((p\Oq)^g,\ +_{F_\tau}\big)$.
\end{theorem}
\begin{proof}
\leanok
For $c \in J^1$: $\redJ(c+\tau) = \redJ(\tau)$, and the reduction of the class of an
integral datum with squarefree reduced $u$-polynomial is the class of the reduced datum
(split $u$ over $L$ with integral roots --- squarefree $\bar u$ lifts to separability
--- factor the Mumford ideal into point ideals on both levels, and match the
reductions). So $\redJ(c+\tau) = [(\bar u_\tau, \bar v_\tau)]$, and
Lemma~\ref{lem:reduce-commutes} shows the reduced representative of $c + \tau$ is
integral with reduction $(\bar u_\tau, \bar v_\tau)$. Its $u$-part is the chart
polynomial at explicit coordinates $z \in (p\Oq)^g$ (the displacement coefficients),
and its $v$-part is the chart square root, by uniqueness of square roots with a
B\'ezout certificate at the common residue datum. Conversely every chart class reduces
to $\redJ(\tau)$, so chart classes translated by $-\tau$ lie in $J^1$; injectivity is
uniqueness of reduced Mumford representatives (Lemma~\ref{lem:mumford-exists}); the
group compatibility is the pointwise class computation of
Theorem~\ref{thm:Ftau-assoc} together with $[(u,-v)] = [(u,v)]^{-1}$.
\end{proof}

\begin{definition}[The logarithm on the kernel of reduction]
\label{def:Lam}
\lean{ChabautyColeman.Lam, ChabautyColeman.Lam_injective, ChabautyColeman.Lam_add,
  ChabautyColeman.ftauFoldFGL, ChabautyColeman.kerScaledCoords,
  ChabautyColeman.FormalGroupLaw.scaledLog}
\leanok
\uses{thm:Ftau-assoc, thm:log-iso, thm:chart-mem}
$\Lam := \log_{F_\tau} \circ\ (\text{chart coordinates of } \cdot + \tau)\colon
J^1(\Qq) \to \Qq^{\,g}$, a group homomorphism by Theorems~\ref{thm:chart-mem}
and~\ref{thm:log-iso}. ($\Lam$ is moreover injective; for $e = 1$ this is formalized in
\texttt{HyperellipticCurves.Archive.LogInjective}, and the main theorems do not depend on it.)
\emph{Formalization.} $\Lam$ is defined in the $\pi$-scaled coordinates of
Theorem~\ref{thm:log-iso}: its value is the total evaluation of the integral $\pi$-scaled
logarithm at $\mathbf y$, where $\pi\mathbf y$ is the chart coordinate vector of
$c + \tau$; additivity combines the group compatibility of
Theorem~\ref{thm:chart-mem} with Theorem~\ref{thm:log-iso}.
\end{definition}

\begin{lemma}[Rationality of the invariant differential]
\label{lem:delta-rational}
\lean{ChabautyColeman.invariantDiffFold_rational, ChabautyColeman.DatumFamRing,
  ChabautyColeman.ftauFamSpec,
  ChabautyColeman.traceDiff_prod, ChabautyColeman.traceDiff_foldPipe,
  ChabautyColeman.sum_traceMat_ftauFoldPipe,
  ChabautyColeman.sum_traceMat_mul_setYZero_pderiv_FtauFold,
  ChabautyColeman.traceMat_map_setYZero_eq_C,
  ChabautyColeman.det_traceMatrix_ne_zero,
  Polynomial.exists_bezout_of_isUnit_resultant,
  ChabautyColeman.ftauFoldPipe_eq_of_certs, ChabautyColeman.FtauFold_eq_of_certs}
\leanok
\uses{def:Ftau, def:fgl-log}
$M(X) = (\partial F_\tau/\partial Y)(X,0)$ is the power-series expansion of a matrix of
\emph{rational} functions of the full Mumford datum, invertible at $0$; hence
$\delta = M^{-1}dX$ is a $g$-tuple of rational $1$-forms on the chart. (Rationality
holds in the coordinates of the full datum; the chart square root is transcendental in
the chart coordinates alone.)
\end{lemma}
\begin{proof}
\leanok
$F_\tau$ is the expansion of a composite of rational maps
(Definitions~\ref{def:cantor-map} and~\ref{def:Ftau}); differentiation preserves
rationality, and $M(0) = I$. \emph{Formalization.} Each entry of $M$ is exhibited as
the specialization, at the chart data with $Y = 0$, of a universal element of a
three-block coefficient ring localized away from an explicit list of denominators, so
that the entries re-specialize along any family of Mumford data. The route is through
the trace form of Theorem~\ref{thm:omega-prime-global}: the invariance identity of the
fold specializes at $Y = 0$ to $T(X)\cdot M(X) = T_2$ with $T$ the trace matrix of the
$X$-side chart and $T_2$ constant, so $M = T^{-1}T_2$ is determined by the starting
data alone, and all four denominators (three block resultants and the determinant of
the trace matrix) specialize to units unconditionally.
\end{proof}

\begin{definition}[Curve-to-chart maps]
\label{def:phiD}
\leanok
\lean{ChabautyColeman.pairV, ChabautyColeman.pair_dvd, ChabautyColeman.diskY_sq,
  ChabautyColeman.movingPair_dvd, ChabautyColeman.exists_isPipelineCerts_of_center,
  ChabautyColeman.exists_isPipelineCerts_powerSeries, ChabautyColeman.chartCoords,
  ChabautyColeman.phiPipeline, ChabautyColeman.phiZeta,
  ChabautyColeman.phiPipeline_map_constantCoeff, ChabautyColeman.constantCoeff_phiZeta}
\uses{def:cantor-map, def:chart, lem:tau-exists, lem:aj-inverse}
For an auxiliary point $Q$ in good position relative to $\tau$, the rational map
$\phi_Q \colon C \dashrightarrow \text{chart}$ sends $P$ to the chart coordinates of
$[P + \iota Q - 2\infty] + \tau$; it is regular at all $P$ avoiding an explicit proper
closed locus. On a residue disk with parameter $z$ and centre $P_0$, the map is
realized as a three-term Cantor computation over $\Oq[[z]]$ on the data of the moving
pair $[P_z + \iota Q]$, the $\tau$-datum, and the constant pair $[Q + \iota P_0]$; its
chart coordinates form the coordinate series $\zeta(z)$. At $z = 0$ the computation
returns the $\tau$-datum exactly --- the two pair classes cancel and reduced
representatives are unique --- so $\zeta(0) = 0$. As everywhere, the B\'ezout data of
the steps are parameters; their existence reduces to residue-level genericity
(Lemma~\ref{lem:global-generic}) by Nakayama lifting. The computation is ordered so
that all support overlaps are of doubling type: cancelling overlaps (a point meeting
its involute) are genuinely singular for the certificate-based composition, which is
what shapes the good-position locus.
\end{definition}

\begin{lemma}[Expansion injectivity]
\label{lem:exp-inj}
\lean{ChabautyColeman.funcExpansionGeneric_eq_zero_iff,
  ChabautyColeman.funcExpansionGeneric_injective}
\leanok
\uses{lem:disk-param}
A rational differential on $C$ over $\Qq$ whose full expansion at one point of
$C(\Qq)$ vanishes is zero; two rational differentials with equal expansions at one
point are equal.
\end{lemma}
\begin{proof}
\leanok
Stated for functions $P + Q\,y$ ($P, Q \in \Qq[x]$) at an affine non-Weierstrass point,
to which rational functions and differentials reduce by clearing denominators: the
expansion times its conjugate expansion is the Taylor series of the norm polynomial
$P^2 - fQ^2$, which is nonzero for $(P,Q) \neq (0,0)$ by the degree-parity count.
\emph{Formalization.} The main development compares series through uniqueness of
reduced representatives instead; this lemma is formalized in
\texttt{HyperellipticCurves.Archive.ExpInj}, and the main theorems do not depend on it.
\end{proof}

\begin{theorem}[The differentials $\omega'_i$]
\label{thm:omega-prime-global}
\lean{ChabautyColeman.omegaPrime, ChabautyColeman.pderiv_subst_log,
  ChabautyColeman.cantorReduce_fst, ChabautyColeman.map_reduceCert,
  ChabautyColeman.exists_reduceCert_map, ChabautyColeman.natDegree_cantorReduce_of_fullDeg,
  ChabautyColeman.isCoprime_two_mul_of_map, ChabautyColeman.isCoprime_cantorStage₂_of_map,
  ChabautyColeman.isCoprime_cantorPipeline_of_map, ChabautyColeman.isCoprime_phiStage₂_of_residue,
  ChabautyColeman.isCoprime_phiPipeline_of_residue, ChabautyColeman.isCoprime_wStage₂_of_residue,
  ChabautyColeman.isCoprime_wPipeline_of_residue,
  ChabautyColeman.subst_log_eq_antideriv, ChabautyColeman.mk0_phiPipeline_eval,
  ChabautyColeman.phiPipeline_eq_of_aux,
  ChabautyColeman.cantorPipeline_eq_cantorPipeline_of_mk0_eq,
  ChabautyColeman.mk0_pair_mul_mk0_pair_chain,
  ChabautyColeman.FractionTopology.valuedFractionRing,
  ChabautyColeman.FractionTopology.continuous_algebraMap,
  ChabautyColeman.FractionTopology.isTopologicallyNilpotent_algebraMap,
  ChabautyColeman.hasSum_algebraMap_evalT_scaledLog, ChabautyColeman.hasSum_algebraMap_Lam,
  ChabautyColeman.redJ_diskPointDiff_eq_one, ChabautyColeman.aj_eq_mk0_mumfordIdeal,
  ChabautyColeman.FractionTopology.isUniformInducing_algebraMap,
  ChabautyColeman.FractionTopology.instCompleteSpace,
  ChabautyColeman.phiPipeline_monic_dvd_degree_lt,
  ChabautyColeman.map_evalRingHom_eq_add_chartCoords,
  ChabautyColeman.map_evalRingHom_eq_chartV_map,
  ChabautyColeman.chartClass_eq_diskPointDiff_mul_tauClass,
  ChabautyColeman.J1_equiv_polydisk_diskPointDiff, ChabautyColeman.summable_logSubstFamily,
  ChabautyColeman.hasSum_coeff_subst_log, ChabautyColeman.hasSum_Lam_antiderivOmegaPrime,
  ChabautyColeman.HypField.deriv, ChabautyColeman.globPipeline, ChabautyColeman.globZeta,
  ChabautyColeman.globOmega, ChabautyColeman.globExpand,
  ChabautyColeman.map_globExpand_globPipeline, ChabautyColeman.globExpand_globZeta,
  ChabautyColeman.globExpand_deriv, ChabautyColeman.globExpand_globOmega,
  ChabautyColeman.regularDiffBasis, ChabautyColeman.map_chartU_eq_of_chartCoords,
  ChabautyColeman.map_chartV_eq_of_chartCoords, ChabautyColeman.map_chartU_eq_phiPipeline_fst,
  ChabautyColeman.map_chartV_eq_phiPipeline_snd, ChabautyColeman.globExpand_comp_datumFamSpec,
  ChabautyColeman.isUnit_datumFamSpec_glob, ChabautyColeman.exists_globN,
  ChabautyColeman.HypRing.exists_toHypField_mul_eq,
  ChabautyColeman.HypRing.exists_toHypField_mul_eq_finite,
  ChabautyColeman.toHypField_eq_of_bezout, ChabautyColeman.HypField.deriv_toHypField_C,
  Polynomial.topCoeffMod, Polynomial.topCoeffMod_eq_zero_of_natDegree_lt,
  Polynomial.topCoeffMod_mul_monic, Polynomial.topCoeffMod_neg,
  Polynomial.topCoeffMod_congr, Derivation.dotDer, Derivation.coeff_dotDer,
  Derivation.dotDer_eq_zero_of_coeff, ChabautyColeman.traceDiff,
  ChabautyColeman.traceDiff_eq_zero_of_dotDer_eq_zero, ChabautyColeman.dvd_sub_of_bezout,
  ChabautyColeman.traceDiff_mul, ChabautyColeman.traceDiff_reduce,
  ChabautyColeman.sub_mul_topCoeffMod_pair, ChabautyColeman.traceDiff_pair,
  Derivation.natDegree_dotDer_le, Polynomial.natDegree_sq_sub_of_odd,
  ChabautyColeman.traceDiff_congr, ChabautyColeman.exists_bezout_of_dvd_sub,
  ChabautyColeman.isCoprime_two_mul_of_dvd_add,
  ChabautyColeman.traceDiff_reduce_of_natDegree_lt, Derivation.dvd_dotDer_of_dvd_sub,
  ChabautyColeman.traceDiff_reduce_of_dvd_dotDer_of_natDegree_lt,
  ChabautyColeman.traceDiff_reduce_of_dvd_dotDer,
  ChabautyColeman.traceDiff_cantorReduce_of_dvd_dotDer,
  ChabautyColeman.traceDiff_cantorPipeline_of_dvd_dotDer, ChabautyColeman.map_residue_eq,
  ChabautyColeman.traceMat_wPipeline, ChabautyColeman.wU₁_map_seriesResidue,
  ChabautyColeman.wV₁_map_seriesResidue, ChabautyColeman.isPipelineAvoid_w_of_residue,
  ChabautyColeman.cantorCompose_eq_of_isCoprime, ChabautyColeman.isCoprime_wU₁_map,
  ChabautyColeman.two_mul_X_mul_traceDiff_wPipeline_map,
  ChabautyColeman.eq_aeval_mul_of_two_mul_eq,
  ChabautyColeman.X_mul_omegaPrime_eq_aeval_tracePoly_mul_derivative,
  ChabautyColeman.map_omegaExpansionWeierstrass_eq_sum_smul_omegaPrime_tracePoly,
  ChabautyColeman.aeval_ne_zero_of_coeff_zero_ne_zero, ChabautyColeman.aeval_inv_eq_zero_iff,
  ChabautyColeman.invAevalFracHom, ChabautyColeman.sq_inv_pow_mul_X,
  ChabautyColeman.inftyExpand, ChabautyColeman.traceDiff_single, ChabautyColeman.IsAddCerts,
  ChabautyColeman.cantorStage₂_monic_dvd_degree_lt, ChabautyColeman.traceDiff_cantorStage₂,
  ChabautyColeman.IsAddCerts.map, ChabautyColeman.mk0_cantorStage₂,
  ChabautyColeman.map_infSeries_eq, ChabautyColeman.traceMat_infU,
  ChabautyColeman.kzDeriv_aeval, ChabautyColeman.kzDeriv_inv,
  ChabautyColeman.two_mul_y_mul_traceDiff_infU_map, ChabautyColeman.infSeries_mul_omegaPrime_kz,
  ChabautyColeman.map_omegaExpansionInf_eq_sum_smul_omegaPrime_tracePoly,
  ChabautyColeman.HypField.derivation, ChabautyColeman.HypField.deriv_smul,
  ChabautyColeman.HypField.deriv_algebraMap_baseField, MvPowerSeries.pderiv,
  RatFunc.derivative_C, RatFunc.derivative_smul, ChabautyColeman.dvd_cantorCompose_sub_left,
  ChabautyColeman.dvd_cantorCompose_sub_right, ChabautyColeman.isCoprime_cantorCompose,
  ChabautyColeman.traceDiff_cantorCompose, ChabautyColeman.dvd_cantorReduce_add,
  ChabautyColeman.traceDiff_cantorReduce, ChabautyColeman.traceDiff_cantorPipeline,
  ChabautyColeman.HypField.deriv_x, Polynomial.topCoeffMod_sum,
  Polynomial.topCoeffMod_C_mul, Derivation.dotDer_C_mul_X_pow, ChabautyColeman.traceMat,
  ChabautyColeman.traceDiff_eq_sum, ChabautyColeman.chartU_map_rename_eq,
  ChabautyColeman.pderiv_coeff_map_C_eq_zero, ChabautyColeman.traceDiff_chartU_map_rename,
  ChabautyColeman.traceDiff_ftauUX_eq_zero, ChabautyColeman.traceDiff_ftauUY,
  ChabautyColeman.IsFoldCerts.isFoldState, ChabautyColeman.traceDiff_prod,
  ChabautyColeman.traceDiff_foldPipe, ChabautyColeman.sum_traceMat_ftauFoldPipe,
  ChabautyColeman.ftauUY_eq_prod, Polynomial.topCoeffMod_map, ChabautyColeman.traceMat_map,
  ChabautyColeman.ftauFoldPipe_map_setYZero,
  ChabautyColeman.sum_traceMat_mul_setYZero_pderiv_FtauFold,
  Polynomial.exists_topCoeffMod_ne_zero, ChabautyColeman.traceMatrix,
  ChabautyColeman.traceMatrix_comm, ChabautyColeman.sum_mul_traceMatrix,
  ChabautyColeman.eq_zero_of_sum_mul_traceMatrix, ChabautyColeman.det_traceMatrix_ne_zero,
  ChabautyColeman.dvd_sub_of_bezout_of_bezout, ChabautyColeman.traceMat_congr,
  ChabautyColeman.traceMat_map_setYZero_eq_C,
  ChabautyColeman.sum_traceMat_mul_setYZero_pderiv_FtauFold_eq_C,
  ChabautyColeman.traceMat_phiPipeline, ChabautyColeman.globExpand_traceMat,
  ChabautyColeman.eq_mul_of_mul_eq,
  ChabautyColeman.constantCoeff_setYZero, ChabautyColeman.residue_eq_of_constantCoeff_sub_mem',
  ChabautyColeman.map_residue_eq', ChabautyColeman.bezout_ftauFoldPipe_map_setYZero,
  MvPowerSeries.substRingHom, MvPowerSeries.substRingHom_apply, MvPowerSeries.substRingHom_X,
  MvPowerSeries.substRingHom_comp_C, MvPowerSeries.constantCoeff_substRingHom,
  MvPowerSeries.map_substRingHom, MvPowerSeries.hasSubst_map_of_constantCoeff_zero,
  ChabautyColeman.map_substRingHom_diffMatrix, ChabautyColeman.ftauFoldFGL_diffMatrix_eq,
  ChabautyColeman.isCoprime_X_sub_C_of_isUnit_eval, ChabautyColeman.isCoprime_pairV,
  ChabautyColeman.ker_constantCoeff_le_jacobson, ChabautyColeman.isCoprime_chartU,
  ChabautyColeman.isCoprime_chartU_map_rename, ChabautyColeman.isCoprime_cantorPipeline,
  ChabautyColeman.HypField.algebraMap_polynomial_injective, ChabautyColeman.HypField.y_ne_zero,
  ChabautyColeman.isCoprime_two_mul_map, ChabautyColeman.isUnit_two_mul_map,
  ChabautyColeman.isCoprime_pair_const, ChabautyColeman.isCoprime_phiU₁,
  ChabautyColeman.isCoprime_phiU₃, ChabautyColeman.isCoprime_neg_map_C,
  ChabautyColeman.isCoprime_ftauFoldPipe, ChabautyColeman.isCoprime_phiPipeline,
  ChabautyColeman.HypRing.toHypField_apply,
  ChabautyColeman.omegaPrime_taylor, ChabautyColeman.isCoprime_of_map,
  ChabautyColeman.ReflectsCoprime, ChabautyColeman.reflectsCoprime_of_ker_le_jacobson,
  ChabautyColeman.reflectsCoprime_of_field, ChabautyColeman.cantorStage₁_map,
  ChabautyColeman.cantorStage₂_map, ChabautyColeman.cantorStage₃_map,
  ChabautyColeman.mvSeriesResidue, ChabautyColeman.mvSeriesResidue_surjective,
  ChabautyColeman.isUnit_of_mvSeriesResidue_ne_zero,
  ChabautyColeman.ker_mvSeriesResidue_le_jacobson, ChabautyColeman.ftauFoldXs_map,
  ChabautyColeman.ftauFoldYs_map, ChabautyColeman.isFoldAvoid_of_map,
  ChabautyColeman.seriesResidue_apply, ChabautyColeman.seriesResidue_C,
  ChabautyColeman.seriesResidue_X, ChabautyColeman.seriesResidue_C_add_X,
  ChabautyColeman.map_C_map_seriesResidue, ChabautyColeman.phiU₁_map_seriesResidue,
  ChabautyColeman.phiV₁_map_seriesResidue, ChabautyColeman.cantorCompose_eq_of_bezout,
  ChabautyColeman.globExpand_algebraMap_base_comp,
  ChabautyColeman.map_globExpand_map_algebraMap, ChabautyColeman.globExpand_x_eq_diskX,
  ChabautyColeman.globExpand_y_eq_diskY, ChabautyColeman.globExpand_globSlope_eq,
  ChabautyColeman.map_globExpand_globU₁, ChabautyColeman.map_globExpand_globV₁,
  ChabautyColeman.map_globExpand_globU₃, ChabautyColeman.map_globExpand_globV₃,
  ChabautyColeman.IsPipelineAvoid, ChabautyColeman.IsFtauFoldAvoid,
  ChabautyColeman.isFtauFoldAvoid_of_center, ChabautyColeman.isPipelineAvoid_phi_of_residue,
  ChabautyColeman.tracePoly, ChabautyColeman.degree_tracePoly_lt,
  ChabautyColeman.two_smul_tracePoly, ChabautyColeman.sum_smul_tracePoly,
  ChabautyColeman.sum_smul_tracePoly_ne_zero, Derivation.leibniz_mul, IsFractionRing.deriv,
  IsFractionRing.derivation, ChabautyColeman.kzDerivation, ChabautyColeman.kzHom,
  ChabautyColeman.IsPipelineAvoid.map, ChabautyColeman.two_mul_y_mul_traceDiff_phiPipeline_map,
  ChabautyColeman.kzOmega, ChabautyColeman.two_mul_kzOmega_eq_sum,
  ChabautyColeman.two_mul_omegaPrime_kz_eq_sum,
  ChabautyColeman.yseries_mul_omegaPrime_eq_taylor_tracePoly_kz,
  ChabautyColeman.map_omegaExpansionGeneric_eq_sum_smul_omegaPrime,
  ChabautyColeman.aeval_eq_zero_iff_of_aeval_eq_X_sq, ChabautyColeman.aevalFracHom,
  ChabautyColeman.weierstrassExpand, ChabautyColeman.weierstrassExpand_toHypField,
  ChabautyColeman.aeval_map_weierstrassSeries,
  ChabautyColeman.map_weierstrassExpand_globPipeline,
  ChabautyColeman.weierstrassExpand_globZeta, ChabautyColeman.derivative_aeval,
  ChabautyColeman.kzDeriv_weierstrassExpand, ChabautyColeman.weierstrassExpand_globOmega,
  ChabautyColeman.X_mul_omegaPrime_eq_aeval_mul_derivative,
  PowerSeries.map_derivative,
  ChabautyColeman.map_omegaExpansionWeierstrass_eq_sum_smul_omegaPrime}
\uses{lem:delta-rational, def:phiD, thm:chart-mem, def:Lam, lem:omega-expansion}
\leanok
Define $\omega'_i := \phi_Q^*\,\delta_i$, rational differentials on $C_{\Qq}$. Then:
\begin{enumerate}
\item[(a)] $\omega'_i$ is independent of the auxiliary point $Q$;
\item[(b)] for every residue disk $D$ with parameter $z$ and all $P_0, P \in D(\Qq)$,
the value $\Lam_i([P]-[P_0])$ --- defined, since same-disk differences lie in $J^1$ ---
is the integral of the expansion of $\omega'_i$ on $D$ from $z(P_0)$ to $z(P)$;
\item[(c)] each $\omega'_i$ is regular, and $(\omega'_1,\dots,\omega'_g)$ is a basis of
$\RegDiff \otimes \Qq$. Explicitly,
\[
  2\,\omega'_l \;=\; \Bigl(\sum_i (A^{-1})_{li}\,x^i\Bigr)\,\frac{dx}{y},
\]
where $A$ is the invertible \emph{trace matrix} of the $\tau$-datum defined in the
proof --- a Hankel matrix depending on $(u_\tau, c_2)$ alone.
\end{enumerate}
The counting argument consumes clause (b), the regularity and explicit form in (c), and
the linear independence of the numerators; clause (a) and the statement that the
$\omega'_i$ span all of $\RegDiff$ are formalized in
\texttt{HyperellipticCurves.Archive.AuxIndep} and
\texttt{HyperellipticCurves.Archive.RegularDiff}, and the main theorems do not depend on them.
\end{theorem}
\begin{proof}
\leanok
\emph{(a)} Two computations of $\phi_Q$ through different auxiliary points produce the
same Mumford datum over $\Oq[[z]]$: the two auxiliary pair classes cancel, and reduced
representatives of a class over $\operatorname{Frac}(\Oq[[z]])$ are unique. Hence
$\zeta$, and with it $\omega'$, does not depend on $Q$.

\emph{(b)} There are three components. \emph{Formal:} by the chain rule and
$d\Lam = \delta$ (Lemma~\ref{lem:fgl-log-diff}), the composite $\Lam_i(\zeta(z))$ is
the antiderivative of $\omega'_i = \sum_j (M^{-1})_{ij}(\zeta(z))\,\zeta_j'(z)$
vanishing at $0$. \emph{Geometric:} the chart coordinates of the same-disk difference
$[P_{z_0}]-[P_0]$, translated by $\tau$, are exactly $\zeta(z_0)$ --- the evaluated
disk computation is the chart datum at $\zeta(z_0)$, its class is
$([P_{z_0}]+[\iota P_0]) + \tau$, and the difference lies in $J^1$ since both points
reduce to the same residue point. \emph{Analytic:} the value of $\Lam$ is a convergent
sum over $\Qq$ (equipped with its valuation topology, for which it is complete), and
the double series obtained by substituting $\zeta(z)$ into the scaled logarithm is
summable, so it may be summed in either order: summing one way gives the value
$\Lam_i(c)$ at the endpoint, summing the other way gives the coefficients of the
antiderivative of $\omega'_i$. Together these prove (b) for a difference against the
disk centre; the general two-point form follows by additivity of $\Lam$.

\emph{(c)} This is proved by a closed computation of the pullback --- the \emph{trace
form} --- with no Riemann--Roch and no case analysis at infinity. For a Mumford datum
$(U,V)$ over a ring with a derivation $D$ acting on coefficients (write $\dot U$ for
the coefficientwise derivative, so $\deg \dot U < \deg U$ as $U$ is monic), set
\[
  \Theta_k(U,V) \;:=\; -\lambda_U\bigl(X^{k-1}\,\dot U\,s\bigr), \qquad
  s := (2V)^{-1} \bmod U, \qquad
  \lambda_u(h) := \mathrm{coeff}_{\deg u - 1}(h \bmod u).
\]
Classically $\Theta_k = \sum_i x_i^{k-1}D(x_i)/(2y_i)$, summed over the points of the
divisor --- the pullback of $x^{k-1}dx/(2y)$ --- but the coefficient description
requires neither the trace form nor separability, and the only inverse it uses,
$(2V)^{-1} \bmod U$, is the B\'ezout cofactor $c_2$ that every step of the computation
carries. Two elementary properties of $\lambda$ --- $\lambda_u(h) = 0$ for
$\deg h \le \deg u - 2$ (the elementary form of Euler's identity
$\mathrm{Tr}(h/u') = 0$), and $\lambda_{uw}(qw) = \lambda_u(q)$ for monic $u, w$, with
no coprimality required (in particular $w = u$ is allowed) --- yield four facts:
\begin{itemize}
\item[(T1)] $\Theta$ is \emph{additive} under Cantor composition, including the
doubling case, where the Chinese remainder theorem fails but the localization property
of $\lambda$ does not;
\item[(T2)] $\Theta$ is \emph{unchanged} by Cantor reduction: with
$R = V^2 - f = U\widetilde U$ one has $\dot R = 2V\dot V$ ($f$ has constant
coefficients), and the two contributions combine into
$\lambda_R(X^{k-1}\dot V) = 0$ by the vanishing property. The degree count is tight and
uses the odd degree of $f$: $2\deg V = 2g+1$ is impossible by parity, and $k \le g$ is
exactly the right range;
\item[(T3)] on a moving pair $\{P, \iota Q\}$, $\Theta_k = x^{k-1}D(x)/(2y)$ directly;
\item[(T4)] data with constant coefficients have $\dot U = 0$, hence $\Theta_k = 0$.
\end{itemize}
Applying (T1)--(T4) along the disk computation of $\phi_Q$ --- one moving block against
two constant ones --- gives, at the generic point,
\[
  y \cdot \Theta_k(U_\zeta, V_\zeta) \;=\; x^{k-1}/2, \qquad k = 1,\dots,g,
\]
a polynomial of degree $k-1 < g$: regularity, outright. Writing
$C_{kj} := -\lambda_{U}(X^{k+j}s)$, a Hankel matrix with only $2g-1$ distinct entries,
one has $\Theta_{k+1} = \sum_j C_{kj}D(\zeta_j)$. Applying (T1), (T2), (T4) in the
formal setting with the derivations $\partial/\partial Y_k$ turns additivity into the
invariance identity $\sum_j C_{ij}(F(X,Y))\,\partial F_j/\partial Y_k = C_{ik}(Y)$, and
setting $Y = 0$ gives $C(X)\,M(X) = C(0)$, i.e.\ $C = A\,M^{-1}$ with $A := C(0)$
\emph{constant}. The matrix $A$ is invertible: it is the Gram matrix, in the monomial
basis of $\Oq[X]/(u_\tau)$, of the residue pairing $(h,q) \mapsto -\lambda_{u_\tau}(hqs)$,
which is nondegenerate whenever $s$ is invertible modulo $u_\tau$ --- guaranteed by the
B\'ezout identity of the $\tau$-datum. Combining, $2\,\varphi_l = \sum_i
(A^{-1})_{li}\,x^i$ for the numerators $\varphi_l$ of the $\omega'_l$: each is a
polynomial of degree $< g$, and, $A$ being invertible, the $\omega'_l$ form a basis.

The identity of (c) is established on each residue disk directly over
$\operatorname{Frac}(K[[z]])$ with the derivation $d/dz$; on a generic disk it reads
$b\,\eta\,\omega'_i = \operatorname{taylor}_a\bigl(\tfrac12\sum_j (A^{-1})_{ij}
x^j\bigr)$, which is the identification consumed by the per-disk count
(Theorem~\ref{thm:disk-bound}). The analogous identities hold on the Weierstrass and
infinity disks; at a Weierstrass centre the second reduction step passes through a
divisor containing a Weierstrass point, and is crossed by a refinement of (T2) --- if
the constant part $B$ of $U = uB$ has $B \mid \dot V$, which holds automatically when
the second block is constant, then $\Theta$ of the reduced datum equals $\Theta$ of the
moving part alone; the orders of $V^2 - f$ and $U$ at the Weierstrass point cancel, so
the reduced data stay Weierstrass-free. At the infinity disk the configuration
$[P_z - \infty] + \tau$ has only two blocks and needs no auxiliary point at all. The
genericity used in this argument --- $\Theta$ must be defined along the computation,
i.e.\ every intermediate divisor must avoid the Weierstrass points --- is exactly the
condition catalogue of Lemma~\ref{lem:global-generic}.
\end{proof}

\begin{remark}[Why a direct holomorphy proof fails]
One might try to prove clause (c) by verifying the holomorphy of $\omega'_i$ place by
place, using denominator-clearing data for two auxiliary configurations together with a
B\'ezout relation between the two denominators. Such data do not exist: by clause (a),
varying the auxiliary point does not move what is being cleared --- the two configurations
compute the same chart coordinates, so the two denominators are denominators of one and the
same function; a B\'ezout relation between them would force the chart coordinates to be
integral over $\Qq[x]$, hence, being also regular at infinity, \emph{constant}, which
happens only in the degenerate case $\zeta \equiv 0$. (Exact computation over the
function field confirms this: for $y^2 = x^5-x+1$, $\tau = [(0,1)]+[(1,1)]-2\infty$,
centre $(-1,1)$, the chart coordinates have denominator $(x^2-x+3)^2$ --- the same for
every auxiliary point.) The trace-form computation in the proof above replaces the
holomorphy verification: it \emph{exhibits} the regular numerator instead of deducing
its existence.
\end{remark}

\begin{remark}[Formalization]
The trace-form differential $\Theta$ is formalized over an arbitrary commutative ring
with a derivation acting on polynomial coefficients (\texttt{HyperellipticCurves.Curve.TraceDiff});
(T2) is proved for domains, where $\deg(V^2 - f) = \max(2\deg V, \deg f)$ --- over a
general ring a nilpotent leading coefficient of $V$ can collapse the degree, and the
identity fails. $\Theta$ does not depend on the choice of B\'ezout cofactor, so each
step of a computation may carry its own certificate and only existence propagates.
Every reduction in the development is a \emph{single} Cantor step: the fold has no
reduction chains, and in the three-block disk computation the outer blocks are pairs,
so each composition has degree $g+2$ and one reduction returns to degree $g$ for every
genus; the B\'ezout data of a computation therefore consist of two composition triples
and two leading-coefficient inverses. The per-disk identities of clause (c) are proved
over $\operatorname{Frac}(K[[z]])$ with the derivation $d/dz$ on all three chart types
(\texttt{HyperellipticCurves.Log.Trace.TraceKz}, \texttt{TraceKzWeierstrass},
\texttt{TraceKzInf}). What separates a chart's trace identity from $\omega'$ --- transport
of the differential matrix along the substitution, inversion of the constant trace matrix,
and the solve --- does not depend on the chart, and is carried out once
(\texttt{two\_mul\_omegaPrime\_kz\_eq\_sum}), the three capstones supplying only their
chart datum and their numerators. The required Weierstrass-avoidance conditions descend to
residue-level conditions on the reduced configuration, since the residue map reflects
coprimality of monic polynomials (Nakayama). Clause (b) rests on the convergence
infrastructure of \texttt{HyperellipticCurves.Series.FractionTopology} (the valuation
topology on $\Qq$ and its completeness, proved via the complete open subring $\Oq$ and
uniformly separated cosets) and the summable double family of
\texttt{HyperellipticCurves.Log.Analytic.ColemanIntegral}.
\end{remark}

\begin{lemma}[Genericity of the global configuration]
\label{lem:global-generic}
\lean{ChabautyColeman.IsPipelineAvoid, ChabautyColeman.isCoprime_two_mul_iff,
  ChabautyColeman.reduceCanon, ChabautyColeman.ReduceGeneric, ChabautyColeman.reduceGeneric_iff,
  ChabautyColeman.ReduceGeneric.reduceCert, ChabautyColeman.ReduceGeneric.fullDeg,
  ChabautyColeman.ReduceGeneric.reduceCanon_natDegree, ChabautyColeman.ReduceCoprimeCanon,
  ChabautyColeman.reduceCoprimeCanon_iff, ChabautyColeman.genericDeg_pair_compose,
  ChabautyColeman.pipelineCanonData, ChabautyColeman.isPipelineCenterCerts_pipelineCanonData,
  ChabautyColeman.isPipelineAvoid_pipelineCanonData,
  ChabautyColeman.isPipelineCenterCerts_of_pair_blocks,
  ChabautyColeman.cantorCompose_snd_eq_add_mul, ChabautyColeman.cantorPipeline_eq_middle,
  ChabautyColeman.isPipelineAvoid_of_pair_blocks, ChabautyColeman.curveVanishing_mul,
  ChabautyColeman.exists_curveTuple_ne_zero, ChabautyColeman.pencilW_spec,
  ChabautyColeman.pencilTangB_ne_zero, ChabautyColeman.finite_setOf_pencil_multiple_root,
  ChabautyColeman.exists_pencil_param, ChabautyColeman.reduceGeneric_pencil_three,
  ChabautyColeman.exists_isPipelineCerts_disk, ChabautyColeman.exists_lift_curve_point,
  ChabautyColeman.pairV_neg, ChabautyColeman.exists_auxiliary_point,
  ChabautyColeman.exists_isPipelineCerts_phi, ChabautyColeman.exists_auxiliary_point_phi,
  ChabautyColeman.eval_pencilW₀_ne_zero_of_deriv,
  ChabautyColeman.pencilPar_ne_zero_of_nonsingular,
  ChabautyColeman.isCoprime_cantorCofactor_pencil,
  ChabautyColeman.movingPair_dvd_weierstrass,
  ChabautyColeman.exists_auxiliary_point_phi_weierstrass, ChabautyColeman.wPipeline,
  ChabautyColeman.wZeta, ChabautyColeman.wPipeline_map_constantCoeff,
  ChabautyColeman.mk0_wPipeline_eval, ChabautyColeman.redJ_diskPointDiff_eq_one_weierstrass,
  ChabautyColeman.chartClass_eq_diskPointDiff_mul_tauClass_weierstrass,
  ChabautyColeman.J1_equiv_polydisk_diskPointDiff_weierstrass,
  ChabautyColeman.hasSum_Lam_antiderivOmegaPrime_weierstrass, ChabautyColeman.diskDiffClassW,
  ChabautyColeman.lambda_glue_weierstrass,
  ChabautyColeman.weierstrassRoot_eq_eval_weierstrassSeries,
  ChabautyColeman.encard_le_of_lambda_eq_zero_weierstrass,
  ChabautyColeman.encard_inter_residueDisk_le_weierstrass,
  ChabautyColeman.encard_inter_residueDisk_le_ordPoint_weierstrass, ChabautyColeman.infC,
  ChabautyColeman.infV, ChabautyColeman.infU, ChabautyColeman.infZeta,
  ChabautyColeman.aeval_reflect_infV, ChabautyColeman.aeval_reflect_infV_sq_sub,
  ChabautyColeman.infV_sq_sub_factor, ChabautyColeman.infU_monic,
  ChabautyColeman.map_constantCoeff_infU, ChabautyColeman.constantCoeff_infZeta,
  ChabautyColeman.exists_infVred, ChabautyColeman.infRoot_eq_eval_infSeries,
  ChabautyColeman.mk0_mumfordIdeal_of_scaled_factor, ChabautyColeman.mk0_infU_eval,
  ChabautyColeman.redJ_diskPointDiff_eq_one_inf, ChabautyColeman.map_constantCoeff_infVred,
  ChabautyColeman.chartClass_eq_diskPointDiff_mul_tauClass_inf,
  ChabautyColeman.J1_equiv_polydisk_diskPointDiff_inf, ChabautyColeman.diskDiffClassInf,
  ChabautyColeman.hasSum_Lam_antiderivOmegaPrime_inf,
  ChabautyColeman.hasSum_lambdaOne_diskDiffClass_inf,
  ChabautyColeman.mapBase_aj_eq_diskDiffClassInf, ChabautyColeman.lambda_glue_inf,
  ChabautyColeman.encard_le_of_lambda_eq_zero_inf,
  ChabautyColeman.encard_inter_residueDisk_le_inf,
  ChabautyColeman.encard_inter_residueDisk_le_ordPoint_inf}
\uses{thm:omega-prime-global, lem:tau-exists, lem:tau-generic, lem:many-points}
\leanok
The genericity hypotheses of Theorem~\ref{thm:omega-prime-global}(c) are satisfiable at
every residue disk: for every nonsingular disk centre there are an auxiliary curve point
$Q = (x_Q, y_Q)$ with its slope datum, and the class $\tau$ of
Lemma~\ref{lem:tau-exists}, such that
\begin{enumerate}
\item $y_Q$ is a unit --- the auxiliary point does not reduce to a Weierstrass point;
\item every intermediate state of the fold computation of $F_\tau$ avoids the
Weierstrass points (this is carried by the datum of Lemma~\ref{lem:tau-exists}, so no
separate condition remains);
\item the reduced per-disk configuration
$[\bar P_0 + \iota\bar Q] + \bar\tau + [\bar Q + \iota \bar P_0]$ avoids the
Weierstrass points at both of its reduced intermediate data.
\end{enumerate}
Like Lemma~\ref{lem:tau-generic}, this is a purely residue-level genericity statement
followed by a Hensel lift.
\end{lemma}
\begin{proof}
\leanok
Condition (2) is Lemma~\ref{lem:tau-generic}, lifted with the datum of
Lemma~\ref{lem:tau-exists}; a single witness configuration discharges both. For a
Mumford datum, $u$ is coprime to $2v$ exactly when it is coprime to $\bar f$, so every
avoidance condition is a condition against the \emph{fixed} curve --- the intermediate
divisor misses the Weierstrass points --- not against the evolving square root.

For conditions (1) and (3), the key observation is that the candidate composed square
roots of the per-disk configuration form an affine \emph{pencil}: with
$N = u_\tau\,(X - a)$ and $V_0$ the interpolant with $V_0 \equiv v_\tau \bmod u_\tau$
and $V_0(a) = b$, they are $V_t = V_0 + t\,N$, and $V_t^2 - f = N\,W_t$ with $W_t$
quadratic in $t$. The auxiliary point is read off the parameter: for any root $x_Q$ of
$W_t$, the pair $(x_Q, -V_t(x_Q))$ is a curve point. Each genericity condition then
becomes ``avoid finitely many $t$'':
\begin{itemize}
\item the degree condition of the first reduction step is $t \neq 0$;
\item transversality at the centre, the condition $u_\tau(x_Q) \neq 0$, and the
coprimality of the reduced intermediate datum with $\bar f$ are each \emph{affine} in
$t$ with nonvanishing coefficient --- respectively $2b$ (at a Weierstrass centre
$b = 0$, and instead the constant term $W_0(a) = -f'(a)/u_\tau(a)$ is nonzero, by
differentiating the pencil identity), $2v_\tau(r)$ at the roots $r$ of $u_\tau$, and
$N(r)$ at the roots $r$ of $\bar f$, all nonzero because $\bar\tau$ is in good
position;
\item the simple-root condition on $W_t$ is finite by elimination: a multiple root $x$
of $W_t$ off the support of $N$ puts the curve point $(x, V_t(x))$ on the line
$A\,y + B = 0$ with $A = 2(V_0'N - V_0N')$ and $B = 2fN' - f'N$, so $x$ is a root of
$A^2f - B^2$; that polynomial is nonzero by degree parity (odd against even), and
$B \neq 0$ because its top coefficient is the integer $2(g+1) - (2g+1) = 1$, nonzero in
\emph{every} characteristic --- no separability of $u_\tau$ and no discriminant is
needed.
\end{itemize}
The third composition produces $-V_t$ on the nose, so its degree condition coincides
with the first; and the full configuration returns its middle datum --- the two pair
blocks are inverse classes --- so the final avoidance condition reduces to the good
position of $\bar\tau$. A good parameter therefore exists over the algebraically closed
residue field, uniformly at every nonsingular centre ($b \neq 0$, or $a$ a simple root
of $\bar f$); at a Weierstrass centre the reduced intermediate data remain
Weierstrass-free because the orders of $V^2 - f$ and of the $u$-part at the centre
cancel. Finally, the residue-level witness lifts by Hensel's lemma: the centre and the
auxiliary point lift as smooth points, the B\'ezout data lift by Nakayama, and all
conditions are open.
\end{proof}

\begin{remark}[Formalization]
The per-disk statement is
\texttt{exists\_isPipelineCerts\_disk} (residue level, all clauses at a nonsingular
centre) and \texttt{exists\_auxiliary\_point} /
\texttt{exists\_auxiliary\_point\_phi} (Hensel-lifted, with the certified computation
over $\Oq[[z]]$), with the Weierstrass-centre forms
\texttt{exists\_auxiliary\_point\_weierstrass} /
\texttt{exists\_auxiliary\_point\_phi\_weierstrass}. The infinity disk needs no
genericity at all: its centre is $\infty$ itself, so the class to expand is
$[P_z - \infty] + \tau$ --- no auxiliary point and no pair blocks; the computation is
one composition and one reduction in closed form, carried in scaled integral
coordinates because the moving point has non-integral coordinates, and the disk
identities (the chart class, the kernel membership, the value of the pairing, and the
per-disk bound of Theorem~\ref{thm:disk-bound}) are established on it directly. On the
Weierstrass and infinity disks the identification of the two disk parametrizations in
use is part of the statement chain. All conditions are produced for every genus
$\ge 1$.
\end{remark}

\begin{definition}[The pairing $\lambda$ on $J^1(\Qp)$, normalized]
\label{def:lambda1}
\leanok
\lean{ChabautyColeman.lambdaOne, ChabautyColeman.lambdaOne_mul,
  ChabautyColeman.lambdaOne_add_left, ChabautyColeman.lambdaOne_smul_left}
\uses{thm:omega-prime-global, lem:global-generic, def:Lam, def:regdiff}
Write $\omega = \sum_i c_i(\omega)\,\omega'_i$ in the basis of
Theorem~\ref{thm:omega-prime-global}, and set
$\lambda_\omega := \sum_i c_i(\omega)\,\Lam_i$ on $J^1(\Qq)$, restricted to $J^1(\Qp)$.
The pairing is linear in $\omega$ and additive in the class. \emph{Formalization:}
$\lambda$ is implemented in $\omega'$-coordinates --- for a coefficient vector $c$ the
functional $\sum_i c_i\,\Lam_i$ on the kernel of reduction; the identification
$c = c(\omega)$ is the explicit trace-form coordinate identity of
Theorem~\ref{thm:omega-prime-global}(c).
\end{definition}

\begin{lemma}[Rationality of values]
\label{lem:values-rational}
\leanok
\lean{ChabautyColeman.lambdaOne_rational_values,
  ChabautyColeman.FractionTopology.continuous_ringEquiv,
  ChabautyColeman.FractionTopology.valuation_ringEquiv}
\uses{def:lambda1, thm:omega-prime-global}
For $\omega \in \RegDiff$ (over $\Qp$) and $c \in J^1(\Qp)$:
$\lambda_\omega(c) \in \Qp$.
\end{lemma}
\begin{proof}
\leanok
$\lambda_\omega$ is characterized Galois-equivariantly by
Theorem~\ref{thm:omega-prime-global}(b) at $\Qp$-rational disks together with
additivity ($\tau$ drops out), so Frobenius-fixed inputs give fixed outputs. The
analytic content is that a continuous field automorphism transports the convergent
characterization of the disk values: any field automorphism of $\Qq$ over $\Qp$ is an
isometry of the valued field, hence continuous, and uniqueness of limits transports the
values along the induced permutation of the disks. \emph{Formalization.} The
annihilator argument of Chapter~5 runs with values in $\Qq$ directly, so the main theorems
do not depend on this lemma; it is formalized in
\texttt{HyperellipticCurves.Archive.RationalValues}.
\end{proof}

\begin{definition}[The pairing $\lambda$ on $J^1(K_0)$]
\label{def:lambda}
\lean{ChabautyColeman.lambda, ChabautyColeman.J1,
  ChabautyColeman.lambda_mul, ChabautyColeman.lambda_add_left,
  ChabautyColeman.lambda_smul_left, ChabautyColeman.lambdaOne_inv,
  ChabautyColeman.mem_J1_iff, ChabautyColeman.mem_J1_of_mapBase_eq,
  ChabautyColeman.mem_J1_of_mapBase_eq_div}
\leanok
\uses{def:lambda1, thm:redJ-hom, lem:values-rational}
Let $J^1(K_0) := \ker\big(\J{K_0} \to \J{K} \to \J{\Fpbar}\big)$ be the group of
classes over the arithmetic base field whose base change reduces trivially, and let
$\lambda_\omega$ on $J^1(K_0)$ be the value of Definition~\ref{def:lambda1} at the base
change of the class. It is additive in the class and linear in $\omega$.
\end{definition}

\begin{remark}[Formalization]
The pairing $\lambda$ is defined on $J^1$ only, and this suffices: every class the
counting argument pairs is a same-disk difference of rational points, which lies in
$J^1$, and the rank of $J^1$ is bounded by the rank of $\J{K_0}$ by flatness of $\QQ$
over $\ZZ$ (Theorem~\ref{thm:annihilator}). So there is no need to extend $\lambda$ to
all of $\J{\Qp}$ by $\lambda(c) := \tfrac1N \lambda(Nc)$, with $N$ the index of $J^1$
(Theorem~\ref{thm:J1-index}), and the finite-index and Frobenius apparatus of
Theorem~\ref{thm:J1-index} does not enter the proof.
\end{remark}

\begin{theorem}[Interface theorem: linearity and the glue identity]
\label{thm:glue}
\lean{ChabautyColeman.lambda_mul, ChabautyColeman.lambda_add_left,
  ChabautyColeman.lambda_smul_left, ChabautyColeman.diskDiffClass,
  ChabautyColeman.hasSum_lambdaOne_diskDiffClass, ChabautyColeman.lambda_glue,
  ChabautyColeman.mapBase_aj_div_eq_diskDiffClass, ChabautyColeman.mapBase_aj,
  ChabautyColeman.HypJac.congr_aj}
\leanok
\uses{def:lambda, thm:omega-prime-global, def:disk-antider}
$\lambda \colon \RegDiff \to \operatorname{Hom}(J^1(\Qp), \Qp)$ is $\Qp$-linear, and
for every residue disk $D$ with parameter $z$ and all $P, P_0 \in C(\Qp) \cap D$ (the
same-disk difference lies in $J^1(\Qp)$, since $\redJ$ kills it):
\[ \lambda_\omega\big(\aj(P) - \aj(P_0)\big) \;=\; G^{\omega}_D(z(P)) - G^{\omega}_D(z(P_0)). \]
\end{theorem}
\begin{proof}
\leanok
Linearity comes with Definitions~\ref{def:lambda1} and~\ref{def:lambda}. For the glue
identity: same-disk differences lie in $J^1(\Qp)$ (Theorem~\ref{thm:redJ-hom}); apply
Theorem~\ref{thm:omega-prime-global}(b) with the basis expansion of $\omega$; the two
antiderivatives agree up to a constant that cancels in the difference. That the base
change of the difference of the two Abel--Jacobi classes is the same-disk difference
used in Theorem~\ref{thm:omega-prime-global}(b) rests on the naturality of the
Abel--Jacobi map in the base field, composed with base change along the extension.
\end{proof}

\begin{remark}[Genus-two shortcut]
\label{rem:g2-cert}
\uses{def:lambda}
For $g = 2$ the constructions of Definitions~\ref{def:Ftau}--\ref{def:lambda} could be
replaced by the explicit Grant--Flynn formal group and local parameters, imported as
polynomial-identity certificates; this would instantiate Theorem~\ref{thm:glue}
independently of Theorems~\ref{thm:Ftau-assoc} and~\ref{thm:omega-prime-global}. The
route is not needed. The script \texttt{prototype/g2\_formal\_group.py} carries out its
first steps in exact arithmetic --- Cantor composition, the chart series, and $F_\tau$ for
$y^2 = x^5 - x + 1$ at $p = 11$ --- and verifies integrality, the unit laws,
commutativity, associativity and $M(0) = I$ up to total degree~$5$.
\end{remark}
